\input{preamble.tex}

\usepackage{tikz-cd}

\DeclareMathOperator{\Char}{char}
\DeclareMathOperator{\im}{im}
\DeclareMathOperator{\coker}{coker}


\title{\huge{MAT224 - Commutative Algebra}}
\author{\LARGE{Thobias Høivik}}
\date{Fall 2026}

\begin{document}
\maketitle

\newpage
\tableofcontents

\newpage
\section*{Intro}
These are my own personal notes in commutative algebra, 
a course I will hopefully be taking in the fall of 2026. 
Note that as these are for myself there might be gaps, inconsistencies
and whatnot which might make it hard to follow so 
keep that in mind if you are considering using my notes for 
yourself. 

\section{Rings}
Throughout the course, and these notes, a ring means 
a commutative ring with unity (meaning multiplicative 
identity).
So every ring $R$ has distinguished elements 
$1_R, 0_R \in R$ and multiplication satisfies 
$ab = ba$. 

\begin{defn}[(Commutative Unital) Ring]
    A ring is a set $A$ together with binary 
    operations $+,\cdot : A \times A \to A$ such that 
    \begin{enumerate}
        \item $(A,+)$ is an abelian group with 
            identity $0$, 
        \item multiplication is associative, 
        \item multiplication is commutative, 
        \item multiplication distributes over addition, 
        \item there exists $1 \in A$ such that 
            $1a = a$ for all $a \in A$. 
    \end{enumerate}
\end{defn}
\begin{ex}
    Examples of familiar rings include 
    \[
        \mathbb Z, \ \mathbb Q, \ \mathbb R, \ \mathbb C
    \]
    with the usual addition and multiplication. 
\end{ex}
If $A$ is a ring, then the polynomial ring $A[x]$ is 
also a ring. More generally, 
$A[x_1,\ldots,x_n]$ is the ring of polynomials in $n$ commuting 
variables with coefficients in $A$. 

\begin{defn}[Ring homomorphism]
    Let $A$ and $B$ be rings. A function 
    \[
        \phi : A \to B
    \]
    is a ring homomorphism if 
    \begin{enumerate}
        \item $\phi(1_A) = 1_B$, 
        \item $\phi(a + a') = \phi(a) + \phi(a')$, 
        \item $\phi(aa') = \phi(a) \phi(a')$. 
    \end{enumerate}
    If $\phi$ is a bijection, then it is 
    called an isomorphism. 
\end{defn}
The condition $\phi(1_A) = 1_B$ is a convention that 
one may or may not choose, but in this course it is 
not optional, and we say all homomorphisms must explicitly 
satisfy this. 

\begin{prop}
    Let $A$ be a ring. There exists a unique ring homomorphism 
    \[
        \iota : \mathbb Z \to \mathbb A.
    \]
    It is given by 
    \[
        \iota(n) = n \cdot 1_A.
    \]
\end{prop}
\begin{proof}
    Since $\iota$ is a rinh homomorphism, it must satisfy 
    \[
        \iota(1) = 1_A. 
    \]
    Then additivity forces 
    \[
        \iota(2) = \iota(1+1) = \iota(1) + \iota(1) 
        = 1_A + 1_A, 
    \]
    and similarly 
    \[
        \iota(n) = \sum_{k=1}^n 1_A = 
        \underbrace{1_A + \cdots + 1_A}_{\text{n times}} 
    \]
    for $n > 0$. Also, 
    \[
        \iota(-n) = -\iota(n).
    \]
    This there is at most one such homomorphism. 
    It is straightforward to check that the function 
    $n \mapsto n \cdot 1_A$ preserves addition, multiplication, 
    and sends $1$ to $1_A$. Hence it is a ring homomorphism. 
\end{proof}

\begin{defn}[Characteristic]
    The characteristic of a ring $A$, denoted $\Char(A)$,
    is the unique integer $n \geq 0$ such that 
    \[
        \ker(\mathbb Z \to A) = (n). 
    \]
\end{defn}
\begin{ex}
    Here are the characteristics of a few familiar 
    rings: 
    \begin{itemize}
        \item $\Char(\mathbb Z) = 0$,  
        \item $\Char(\mathbb F_p) = p$, 
        \item $\Char(\mathbb Q) = 0$. 
    \end{itemize}
\end{ex}

Intuitively, the characteristic is the smallest 
$n \in \mathbb Z^+$ such that $n \cdot 1_A = 0_A$ with 
$n := 0$ if no such $n$ exists, and in fact, the characteristic 
is usually defined as such in a first course in abstract 
algebra. 

\begin{prop}[Universal property of $\mathbb Z$-polynomials]
    Let $A$ be a ring and let $a \in A$. There exists 
    a unique ring homomorphism 
    \[
        \phi_a : \mathbb Z[x] \to A
    \]
    such that 
    \[
        \phi_a(x) = a. 
    \]
    It is given by 
    \[
        \phi_a\left(\sum_{k=0}^n c_kx^k\right) = c_0 \cdot 1_A 
        + c_1 a + \cdots + c_n a^n.
    \]
\end{prop}
\begin{proof}
    A ring homomorphism $\phi : \mathbb Z[x] \to A$ must agree 
    with the unique homomorphism $\mathbb Z \to A$ on constant 
    polynomials. 

    If $\phi(x) = a$, then multiplicativity forces 
    \[
        \phi(x^n) = a^n. 
    \]
    Additivity then forces 
    \[
        \phi(c_0 + c_1x + \cdots + c_nx^n) 
        = c_0 \cdot 1_A + c_1a + \cdots + c_na^n. 
    \]
    Thus there is at most one such homomorphism. 

    Conversely, the formula above defines a ring homomorphism, by 
        the usual rules for polynomial addition and multiplication.
\end{proof}
So homomorphism $\mathbb Z[x] \to A$ are the same thing as 
choices of an element $a \in A$. Intuitively, $\mathbb Z[x]$ is 
the "freest" ring generated by one element. 

\begin{defn}[Ideal]
    Let $A$ be a ring. An ideal $I \subseteq A$ is a subset 
    such that: 
    \begin{enumerate}
        \item $I$ is an additive subgroup of $A$, 
        \item for every $a \in A$ and every $x \in I$, we have 
            \[
                ax \in I. 
            \]
    \end{enumerate}
\end{defn}

\begin{defn}[Principal ideal]
    Let $A$ be a ring and $a \in A$. The principal ideal 
    by $a$ is 
    \[
        (a) = \{ra \bigm| r \in A\}. 
    \]
    An ideal is called principal if it is of the form 
    $(a)$ for some $a \in A$.
\end{defn}
More generally, if $a_1, \ldots, a_n \in A$, then 
\[
    (a_1,\ldots,a_n) = 
    \{r_1a_1 + \cdots + r_na_n \bigm| r_i \in A\}.
\]
This is the smallest ideal containing all of $a_1,\ldots,a_n$. 
\begin{ex}
    $(x,y) \subseteq k[x,y]$ is such an ideal. 
    \[
        (x,y) = 
        \{f(x,y)x + g(x,y)y \bigm| f,g \in k[x,y]\}. 
    \]
    It is not the whole ring. It consists precisely of 
    polynomials with zero constant term: 
    \[
        (x,y) = \{h \in k[x,y] \bigm| h(0,0) = 0\}.
    \]
\end{ex}

\begin{ex}
    $(0) = \{0\}$ is a (very uninteresting) ideal. 
    $(1) = A$ is another (very uninteresting) ideal. 
    For $n \in \mathbb Z$ we have that $(n) = n\mathbb Z$ is 
    an ideal of $\mathbb Z$. 
\end{ex}

\begin{prop}
    Let $\phi : A \to B$ be a ring homomorphism. Then 
    \[
        \ker \phi = \{a \in A \bigm| \phi(a) = 0\}
    \]
    is an ideal of $A$.
\end{prop}

\begin{proof}
    First we show that $\ker \phi$ is an additive subgroup. 
    Since $\phi(0) = 0$, we have $0 \in \ker \phi$. 
    If $x,y \in \ker \phi$, then 
    \[
        \phi(x-y) = \phi(x) - \phi(y) = 0 - 0 = 0. 
    \]
    Hence $x-y \in \ker \phi$. 
    Now let $a \in A$ and $x \in \ker \phi$. Then 
    \[
        \phi(ax) = \phi(a)\phi(x) = \phi(a) \cdot 0 = 0. 
    \]
    Thus $ax = \ker \phi$. Therefore $\ker \phi$ is an ideal. 
\end{proof}

\begin{defn}[Quotient ring]
    Let $I \subseteq A$ be an ideal. The quotient ring 
    $A/I$ is the set of additive cosets 
    \[
        a + I
    \]
    with addition and multiplication defined by 
    \[
        (a+I) + (b+I) = (a+b) + I
    \]
    and 
    \[
        (a+I)(b+I) = ab + I. 
    \]
\end{defn}
\begin{prop}
    If $I \subseteq A$ is an ideal, then $A/I$ is a ring. 
\end{prop}
\begin{proof}
    The quotient $A/I$ is already an abelian group under addition because $I$ is an additive subgroup.

    We need to check multiplication is well-defined.

    Suppose

    $$a+I=a'+I$$

    and

    $$b+I=b'+I.$$

    Then

    $$a-a'\in I,\qquad b-b'\in I.$$

    We want to prove

    $$ab+I=a'b'+I,$$

    i.e.

    $$ab-a'b'\in I.$$

    Compute:

    $$ab-a'b'=ab-a'b+a'b-a'b'=(a-a')b+a'(b-b').$$

    Since $a-a'\in I$, and $I$ absorbs multiplication by elements of $A$, we have

    $$(a-a')b\in I.$$

    Similarly,

    $$a'(b-b')\in I.$$

    Since $I$ is closed under addition,

    $$ab-a'b'\in I.$$

    Therefore multiplication is well-defined.

    The ring axioms for $A/I$ then follow from the ring axioms in $A$.
\end{proof}
The quotient map is 
\[
    \pi : A \to A/I, \ a \mapsto a + I.
\]
Its kernel is exactly $I$: 
\[
    \ker \pi = I. 
\]

\begin{thm}[Fundamental homomorphism theorem]
    Let 
    \[
        \phi : A \to B
    \]
    be a ring homomorphism. Then 
    \[
        \im \phi
    \]
    is a subring of $B$, and there is an isomorphism 
    \[
        A / \ker \phi \cong \im \phi 
    \]
    given by 
    \[
        a + \ker \phi \mapsto \phi(a). 
    \]
\end{thm}
\begin{proof}
    Define 
    \[
        \overline{\phi} : A/\ker \phi \to \im \phi
    \]
    by 
    \[
        \overline \phi(a + \ker \phi) = \phi(a). 
    \]

    We first check well-definedness.

    Suppose

    $$a+\ker\varphi=b+\ker\varphi.$$

    Then

    $$a-b\in\ker\varphi.$$

    Thus

    $$\varphi(a-b)=0.$$

    So

    $$\varphi(a)-\varphi(b)=0,$$

    and hence

    $$\varphi(a)=\varphi(b).$$

    Therefore $\overline{\varphi}$ is well-defined.

    It is a ring homomorphism because $\varphi$ is.

    It is surjective onto $\operatorname{im}\varphi$ by definition.

    Finally, its kernel is trivial. Indeed,

    $$\overline{\varphi}(a+\ker\varphi)=0$$

    if and only if

    $$\varphi(a)=0,$$

    if and only if

    $$a\in\ker\varphi,$$

    if and only if

    $$a+\ker\varphi=0+\ker\varphi.$$

    Thus $\overline{\varphi}$ is injective. 
    Hence it is an isomorphism.
\end{proof}
\begin{ex}
    \[
        \mathbb Z[x]/(x^2 + 1) \cong \mathbb Z(i) = 
        \{a + bi \bigm| a,b \in \mathbb Z\}
    \]
    by $f \mapsto f(i)$. 

    We also have 
    \[
        k[x,y]/(x) \cong k(y) 
    \]
    by $f \mapsto f(0,y)$. 

    Lastly 
    \[
        k[x,y] / (x,y) \cong k
    \]
    by $f \mapsto f(0,0)$.
\end{ex}

\begin{thm}
    Let $I \subseteq A$ be an ideal. There is a bijection 
    \[
        \{\text{ideals of} \ A/I\} 
        \leftrightarrow \{\text{ideals} \ J 
        \subseteq A \ \text{such that} \ I \subseteq J\}. 
    \]
    The maps are 
    \[
        K \subseteq A/I \mapsto \pi^{-1}(K) \subseteq A, 
    \]
    and 
    \[
        J \subseteq A, I \subseteq J \mapsto 
        J/I \subseteq A/I. 
    \]
\end{thm}

\begin{proof}
    Let 
    \[
        \pi : A \to A/I
    \]
    be the quotient map. First, if $K \subseteq A/I$ is an ideal, 
    then $\pi^{-1}(K)$ is an ideal of $A$. Since $0 \in K$, and 
    $\ker \pi = I$, we have 
    \[
        I \subseteq \pi^{-1}(K).
    \]
    Conversely, suppose $J \subseteq A$ is an ideal with 
    $I \subseteq J$. Then 
    \[
        J/I = \{j + I \bigm | j \in J\}
    \]
    is an ideal of $A/I$. Now check that these constructions are 
    inverse. If $K \subseteq A/I$, then 
    \[
        \pi(\pi^{-1}(K)) = K
    \]
    because $\pi$ is surjective. If $J \subseteq A$ contains $I$, 
    then 
    \[
        \pi^{-1}(J/I) = J. 
    \]
    Indeed, $a \in \pi^{-1}(J/I)$ if and only if 
    \[
        a + I \in J/I, 
    \]
    which happens if and only if there exists $j \in J$ such that 
    \[
        a + I = j + I.
    \]
    This means 
    \[
        a - j \in I. 
    \]
    Since $I \subseteq J$ and $j \in J$, we get 
    $a \in J$. Hence 
    \[
        \pi^{-1}(J/I) = J. 
    \]
    Therefore the two maps are inverse bijections. 
\end{proof}
\begin{ex}
    The ideals of $\mathbb Z / 12 \mathbb Z$ correspond 
    to ideals of $\mathbb Z$ containing $(12)$. 
    The ideals of $\mathbb Z$ are precisely 
    $(n)$ for $n \geq 0$. We need 
    $(12) \subseteq (n)$. This happens precisely when 
    $n \bigm| 12$. So the ideals of $\mathbb Z/12\mathbb Z$ are 
    \[
        (1), (2), (3), (4), (6), (12) = (0). 
    \]
\end{ex}

\begin{prop}
    Let $k$ be a field. The ideal 
    \[
        (x,y) \subseteq k[x,y]
    \]
    is not principal. 
\end{prop}
\begin{proof}
    Suppose for a contradiction that 
    \[
        (x,y) = (f)
    \]
    for some $f \in k[x,y]$. 
    Since 
    \[
        x \in (x,y) = (f),
    \]
    we have $f \bigm| x$. Since 
    \[
        y \in (x,y) = (f), 
    \]
    we have $f \bigm| y$. However, $x$ and $y$ have no common 
    non-unit factor in $k[x,y]$. Hence $f$ must be a unit. If 
    $f$ is a unit, then 
    \[
        (f) = (1) = k[x,y].
    \]
    But this is impossible, because 
    \[
        1 \not\in (x,y). 
    \]
    Therefore $(x,y)$ is not principal. 
\end{proof}
Thus $k[x]$ is a principal ideal domain when $k$ is a field, 
but $k[x,y]$ is not. 

\subsection*{Exercises}

\begin{prob}
    Prove that if $A$ is a ring with $1 = 0$, then 
    $A$ is the trivial ring. 
\end{prob}
\begin{proof}
    Let $a \in A$. Then 
    \begin{align*}
        a &= a \\ 
        1 \cdot a &= 0 \cdot a \\ 
        a &= 0. 
    \end{align*}
    Thus 
    \[
        A = \{0\}.
    \]
\end{proof}

\begin{prob}
    Let $A$ be a ring. Prove that there is a unique 
    homomorphism $\mathbb Z \to A$. 
\end{prob}
\begin{proof}
    Let $A$ be a ring and 
    let $\phi, \psi : \mathbb Z \to A$ be homomorphisms. 

    Then we require 
    \[
        \phi(1) = \psi(1) = 1_A. 
    \]
    Therefore, for $n > 0$, we have 
    \[
        \phi(n) = \phi\left(\sum_{k=1}^n 1\right) 
        = \sum_{k=1}^n 1_A 
        = \psi\left(\sum_{k=1}^n 1\right) 
        = \psi(n). 
    \]
    For $n < 0$ we have $n = -m$ for some 
    $m > 0$. Then 
    \[
        \phi(n) = \phi(-m) = -\phi(m) 
        = - \sum_{k=1}^n 1_A 
        = -\psi(m) = \psi(-m) = \psi(n).
    \]
    Lastly, 
    \[
        \phi(0) = 0 = \psi(0). 
    \]
    Thus $\phi$ and $\psi$ agree on all 
    of $\mathbb Z$ and hence 
    \[
        \psi = \phi. 
    \]
\end{proof}

\begin{prob}
    Identify $k[x,y] / (x,y-1)$. 
\end{prob}
\begin{proof}[Solution]
    Let $\phi : k[x,y] \to k$ by 
    \[
        \phi(f) = f(0,1). 
    \]
    Then $\phi$ is a ring homomorphism and 
    furthermore it is surjective since, for 
    $a \in k$ we have $f(x,y) := a \mapsto a$. 

    Now we show that 
    $\ker \phi = (x,y-1)$. First of all 
    $x, y-1 \in \ker \phi$ since $x(0,1) = 0$ and 
    $y-1(0,1) = 0$. Thus 
    \[
        (x,y-1) \subseteq \ker \phi.
    \]
    Now suppose $f \in \ker \phi$. Then 
    $f(0,1) = 0$ so $f(x,y) = xg(x,y) + (y-1)h(x,y) + 
    f(0,1)$. Since $f(0,1) = 0$ we have $f(x,y) = xg(x,y) + 
    (y-1)h(x,y)$ so $f(x,y) \in (x, y-1)$. Thus 
    $\ker \phi \subseteq (x,y-1)$. The two together give 
    $(x,y-1) = \ker \phi$. 
    Thus, by the fundamental theorem of ring homomorphisms 
    (or first isomorphism theorem, if you prefer), we have 
    \[
        \frac{k[x,y]}{(x,y-1)} \cong 
        k.
    \]
\end{proof}

\begin{prob}
    List all ideals of $\mathbb Z/18\mathbb Z$. 
\end{prob}
\begin{proof}[Solution]
    The ideals are 
    \begin{itemize}
        \item $(0)$, 
        \item $(1)$,
        \item $(2)$, 
        \item $(3)$, 
        \item $(6)$, 
        \item $(9)$,  
    \end{itemize}
    i.e. $(n)$ where $n \bigm| 18$. 
    Note that $(0) = (18)$ and $18$ divides $18$.
\end{proof}

\section{Special ideals and rings}

\begin{defn}
    Let $A$ be a ring. 

    An element $u \in A$ is a \textbf{unit} if there 
    exists $v \in A$ such that 
    \[
        uv = 1.
    \]
    An element $x \in A$ is a \textbf{zero-divisor} 
    if there exists a nonzero element $y \in A$ 
    such that 
    \[
        xy = 0.
    \]
    A ring $A$ is an \textbf{integral domain} if 
    \[
        A \neq 0
    \]
    and $0$ is the only zero-divisor. 

    Equivalently, $A$ is an integral domain if 
    \[
        ab = 0 \Rightarrow a = 0 \ \text{or} \ b = 0. 
    \]
    A ring $A$ is a \textbf{field} if 
    \[
        A \neq 0
    \]
    and every nonzero element of $A$ is a unit. 
\end{defn}
\begin{ex}
    $\mathbb Z$ is an integral domain, but not a field. 
    The rings 
    \[
        \mathbb Q, \ \mathbb R, \ \mathbb C
    \]
    are fields. 
    The ring $\mathbb Z/4 \mathbb Z$ is not an integral 
    domain. 
\end{ex}

\begin{prop}
    Let $A$ be a ring and $x \in A$. Then 
    \[
        x \ \text{is a unit}
    \]
    if and only if 
    \[
        (x) = (1).
    \]
\end{prop}
\begin{proof}
    Suppose first $x$ is a unit. Then $\exists y \in A$ 
    s.t. 
    \[
        xy = 1. 
    \]
    Therefore $1 \in (x)$. Since every ideal 
    containing $1$ is the whole ring, we get 
    \[
        (x) = A = (1).
    \]
    Conversely, suppose $(x) = (1)$. 
    Then 
    \[
        1 \in (x). 
    \]
    By definition of the principal ideal $(x)$, there exists 
    $y \in A$ such that 
    \[
        1 = yx. 
    \]
    Since $A$ is commutative, 
    \[
        xy = 1.
    \]
    Therefore $x$ is a unit.
\end{proof}

\begin{thm}
    Let $A$ be a ring. Then $A$ is a field if and only if 
    $A$ has exactly two ideals: 
    \[
        (0) \neq (1). 
    \]
\end{thm}
\begin{proof}
    Suppose $A$ is a field. 

    Let $I \subseteq A$ be a nonzero ideal. Choose 
    \[
        0 \neq x \in I. 
    \]
    Since $A$ is a field, $x$ is a unit. Then 
    \[
        (x) = (1). 
    \]
    But 
    \[
        (x) \subseteq I. 
    \]
    Therefore 
    \[
        (1) = A \subseteq I
    \]
    so 
    \[
        I = (1). 
    \]
    So the only ideals are 
    \[
        (0), \ (1). 
    \]

    Conversely, suppose $A$ has two ideals 
    \[
        (0) \neq (1). 
    \]
    Let $0 \neq x \in A$. Then 
    \[
        (x) \neq (0). 
    \]
    Since the only nonzero ideal is $(1)$, we have 
    \[
        (x) = (1). 
    \]
    Then $x$ is a unit and thus since every nonzero element of 
    $A$ is a unit, $A$ is a field. 
\end{proof}

\begin{defn}[Prime ideal]
    Let $A$ be a ring and let $I\subseteq A$ be an ideal. 
    We say $I$ is a \textbf{prime ideal} if

    $$I\neq(1)$$

    and whenever

    $$ab\in I,$$

    we have

    $$a\in I\quad\text{or}\quad b\in I.$$

    Equivalently,

    $$a,b\notin I\implies ab\notin I.$$

    The set of all prime ideals of $A$ is denoted

    $$\operatorname{Spec}(A).$$
\end{defn}

\begin{ex}
    The prime ideals of $\mathbb Z$ are 
    $(0)$ and $(p)$ where $p$ is a prime number. 
\end{ex}

\begin{defn}[Maximal ideal]
    Let $A$ be a ring and let $I\subseteq A$ be an ideal. 
    We say $I$ is a \textbf{maximal ideal} if

    $$I\neq(1)$$

    and the only ideals $J$ satisfying

    $$I\subseteq J\subseteq A$$

    are

    $$J=I$$

    and

    $$J=A.$$

    Equivalently, there is no ideal strictly between $I$ and $A$.

    The set of maximal ideals of $A$ is denoted

    $$\operatorname{Specm}(A).$$
\end{defn}
\begin{ex}
    The maximal ideals of $\mathbb Z$ are $(p)$. 
\end{ex} 

\begin{thm}
    Let $A$ be a ring and let $I \subseteq A$ be an ideal. 
    Then 
    \[
        I \ \text{is prime}
    \]
    if and only if 
    \[
        A/I \ \text{is an integral domain}.
    \]
    Also, 
    \[
        I \ \text{is maximal}
    \]
    if and only if 
    \[
        A/I \ \text{is a field}. 
    \]
\end{thm}

\subsection*{Exercises}

\begin{prob}
    Decide whether $(2) \subseteq \mathbb Z/6\mathbb Z$ 
    is prime, maximal, both, or neither. 
\end{prob}
\begin{proof}[Solution]
    It is maximal. 
    Let $\phi : \mathbb Z/6\mathbb Z \to \mathbb Z/2\mathbb Z$ 
    by 
    \[
        \phi(a) = a \mod 2. 
    \]
    Then this map is surjective and 
    $\ker \phi = (2)$. By the first isomorphism theorem we 
    have 
    \[
        \mathbb Z_6/\ker \phi \cong \mathbb Z_2
    \]
    which is a field and thus $\ker \phi$ is maximal.  
\end{proof}

\begin{prob}
    Decide whether $(x-1) \subseteq k[x]$ is prime or maximal. 
\end{prob}
\begin{proof}[Solution]
    Let $\phi : k[x] \to k$ by 
    \[
        \phi(f) = f(1). 
    \]
    Then $\ker \phi = (x-1)$, the map is surjective so 
    it follows as from before that $(x-1)$ is maximal. 
\end{proof}

\section{Krull's theorem, PIDs/UFDs, local rings}

\begin{thm}[Krull's theorem]
    Let $A$ be a nonzero ring. Then $A$ has a maximal 
    ideal. 

    Equivalently: 
    \[
        A \neq 0 
        \Rightarrow \operatorname{Specm}(A) \neq \emptyset. 
    \]
\end{thm}
This is not at all obvious for arbitrary rings, because 
rings may have infinitely many ideals and no obvious "largest" 
proper ideal. To prove Krull's theorem we will need 
the following set-theoretic principle. 

\begin{lem}[Zorn's lemma]
    Let $S$ be a partially ordered set. Suppose 
    every chain in $S$ has an upper bound in $S$. 
    Then $S$ has a maximal element. 
\end{lem}

Here: 
\begin{itemize}
    \item a \textbf{chain} is a subset where any 
        two elements are comparable, 
    \item an \textbf{upper bound} for a chian 
        $C \subseteq S$ is an element $u \in S$ such 
        that $c \leq u$ for all $u \in C$, 
    \item a \textbf{maximal element} is an element 
        $m \in S$ such that there is no $s \in S$ 
        with $m < s$. 
\end{itemize}

\begin{proof}[Proof of Krull's theorem]
    Let $A \neq 0$. Consider the set 
    \[
        S = \{I \subseteq A \bigm| I \ \text{is an ideal and} \ 
        I \neq A\}.
    \]
    So $S$ is the set of all proper ideals of $A$. 
    We partially order $S$ by inclusion: 
    \[
        I \leq J \Leftrightarrow I \subseteq J. 
    \]
    We want to apply Zorn's lemma. 

    Let $\{I_\lambda\}_{\lambda \in \Lambda}$ be a chain in 
    $S$. Define 
    \[
        I = \bigcup_{\lambda \in \Lambda} I_\lambda.
    \]
    We claim $I \in S$, meaning $I$ is a proper ideal. 

    First, $I$ is an ideal. 

    Let $a,b \in I$. Then there exists $\lambda, \mu \in \Lambda$ 
    such that 
    \[
        a \in I_\lambda, \ b \in I_\mu. 
    \]
    Since $I_\mu$ is an ideal, 
    \[
        a - b \in I_\mu \subseteq I. 
    \]
    Thus $I$ is an additive subgroup. 
    Now let $r \in A$ and $a \in I$. Then $a \in I_\lambda$ for 
    some $\lambda$, so 
    \[
        ra \in I_\lambda \subseteq I. 
    \]
    Thus $I$ absorbs multiplication from $A$ and is therefore 
    an ideal. Second $I \neq A$. 
    If $I = A$, then in particular 
    \[
        1 \in I.
    \]
    So $1 \in I_\lambda$ for some $\lambda$, but then 
    \[
        I_\lambda = A,
    \]
    contradicting $I_\lambda \in S$. Hence $I \neq A$. 

    So $I \in S$, and $I$ is an upper bound for the chain. 
    Therefore every chain in $S$ has an upper bound. 
    By Zorn's lemma, $S$ has a maximal element $m$. 
    This $m$ is a proper ideal of $A$, and no proper ideal 
    strictly contains it. Hence $m$ is a maximal ideal. 

    Thus every nonzero ring has a maximal ideal. 
\end{proof}

\begin{cor}
    Let $I \subseteq A$ be a proper ideal. 
    Then there exists a maximal ideal $m \subseteq A$ such 
    that 
    \[
        I \subseteq m. 
    \]
\end{cor}
\begin{proof}
    Since $I \neq A$, the quotient 
    \[
        A/I
    \]
    is nonzero. By Krull's theorem, $A/I$ has a maximal ideal 
    $\mathfrak m$. 
    Let 
    \[
        \pi : A \to A/I   
    \]
    be the quotient map. Then 
    \[
        m = \pi^{-1}(\mathfrak m)
    \]
    is a maximal ideal of $A$ containing $I$, by 
    the correspondence theorem. 
\end{proof}
\begin{cor}
    Let $A$ be a ring and $x \in A$. Then 
    \[
        x \ \text{is a unit}
    \]
    if and only if 
    \[
        x \not\in m
    \]
    for every maximal ideal $m \subseteq A$. 

    Equivalently: 
    \[
        x \ \text{is not a unit} 
        \Leftrightarrow x \in m \ \text{for some maximal ideal} \ 
        m. 
    \]
\end{cor}
\begin{proof}
    We have: 
    \[
        x \ \text{is not a unit} \Leftrightarrow 
        (x) \neq A. 
    \]
    By the previous Corollary, 
    \[
        (x) \neq A \Leftrightarrow (x) \subseteq m
    \]
    for some maximal ideal $m$. 
    But 
    \[
        (x) \subseteq m \Leftrightarrow x \in m. 
    \]
    Therefore 
    \[
        x \ \text{is not a unit} 
        \Leftrightarrow x \in m
    \]
    for some maximal ideal $m$. 
\end{proof}

\begin{defn}[PID]
    An integral domain $A$ is a 
    \textbf{principal ideal domain}, or PID, if every ideal 
    of $A$ is principal. 
\end{defn}
\begin{ex}
    Some examples of PIDs include 
    \begin{itemize}
        \item $k[x]$ for any field $k$, 
        \item any field $k$, 
        \item $\mathbb Z$. 
    \end{itemize}
\end{ex}

\begin{defn}
    Let $A$ be an integral domain. Let $f \in A$ be neither 
    $0$ nor a unit. We say $f$ is \textbf{irreducible} 
    if whenever 
    \[
        f = gh, 
    \]
    then either $g$ is a unit or $h$ is a unit, 
    meaning $f$ cannot be factored nontrivially. 
\end{defn}
\begin{defn}[UFD]
    An integral domain $A$ is a \textbf{unique factorization domain},
    or UFD, if every nonzero nonunit $f \in A$ can be written as 
    a finite product of irreducibles, 
    \[
        f = p_1 \cdots p_n,
    \]
    and this factorization is unique up to reordering and  
    multiplication by units. 
\end{defn}
\begin{ex}
    Some examples include: 
    \begin{itemize}
        \item $\mathbb Z$, 
        \item $k[x]$ for any field $k$, 
        \item $k[x_1,\ldots,x_n]$ for any field $k$. 
    \end{itemize}
\end{ex}

We remark that PID $\Rightarrow$ UFD and that if 
$A$ is a UFD then $A[x]$ is a UFD. 

\begin{defn}[Local ring]
    A ring $A$ is \textbf{local} if it has exactly 
    one maximal ideal. 
    That maximal ideal is usually denoted 
    \[
        \mathfrak m. 
    \]
    So a local ring is a pair $(A, \mathfrak m)$, 
    where $\mathfrak m$ is the unique maximal ideal. 
\end{defn}

\begin{ex}
    Every field $k$ is local since the only maximal ideal 
    is $(0)$. 

    $\mathbb Z$ is a nonexample as it has maximal ideals 
    $(p)$ for every prime $p$. 
\end{ex}

\begin{prop}
    Let $A$ be a local ring with unique maximal ideal 
    $\mathfrak m$. Then the set of units in $A$ is 
    \[
        A \setminus \mathfrak m. 
    \]
    Equivalently: 
    \[
        x \in \mathfrak m 
        \Leftrightarrow x \ \text{is not a unit}. 
    \]
\end{prop}
\begin{proof}
    By the corollary to Krull's theorem, an element $x \in A$ 
    is a unit if and only if it is contained in no maximal ideal. 
    But $A$ has exactly one maximal ideal, namely 
    $\mathfrak m$. 
    Therefore 
    \[
        x \ \text{is a unit} \Leftrightarrow 
        x \not\in \mathfrak m. 
    \]
    So the units are exactly 
    \[
        A \setminus \mathfrak m. 
    \]
\end{proof}

\begin{prop}
    Let $A$ be a ring. Suppose the set of nonunits 
    \[
        \mathfrak m = \{x \in A \bigm| x \ \text{is not 
        unit}\}
    \]
    is an ideal. 
    Then $A$ is local, and $\mathfrak m$ is its unique 
    maximal ideal. 
\end{prop}
\begin{proof}
    Let $I \subseteq A$ be a proper ideal. 
    We claim 
    \[
        I \subseteq \mathfrak m. 
    \]
    Indeed, if $x \in I$ were a unit, then 
    \[
        1 \in I,
    \]
    because 
    \[
        x^{-1}x = 1. 
    \]
    Then $I = A$, contradicting that $I$ is proper. 
    So every element of $I$ is a nonunit, hence 
    \[
        I \subseteq \mathfrak m. 
    \]
    Thus every proper ideal of $A$ is contained in 
    $\mathfrak m$. 
    Since $\mathfrak m$ is itself a proper ideal, 
    it is the unique maximal ideal. Therefore $A$ is local. 
\end{proof}

\subsection*{Exercises}

\begin{prob}
    \ 
    \begin{enumerate}[label=(\alph*)]
        \item Let $A$ be a ring, and let 
            $f = a_nx^n + \cdots + a_0 \in A[x]$. Prove 
            that if $f$ is a unit of $A[x]$, then 
            $a_0$ is a unit of $A$.
        \item Show that if $a$ is a unit of $A$, then the 
            constant polynomial $a \in A[x]$ is a unit of 
            $A[x]$. 
        \item Show that if $A$ is an integral domain, 
            then the units of $A[x]$ are exactly the constant 
            polynomials $a$ with $a$ a unit of $A$. 
    \end{enumerate}
\end{prob}
\begin{proof}[Proof of (a)]
    Suppose 
    \[
        f = \sum_{k=0}^n a_kx^k
    \]
    is a unit of $A[x]$. Then there exists 
    \[
        g = \sum_{k=0}^m b_kx^k
    \]
    such that 
    \[
        fg = 1. 
    \]
    Let $\phi : A[x] \to A$ by 
    \[
        \phi(f) = f(0).
    \]
    Then $\phi$ is a ring homomorphism, and thus 
    \[
        1_A = \phi(1_{A[x]}) 
        = \phi(fg) = a_0 b_0 
    \]
    so $a_0$ is a unit in $A$. 
\end{proof}
\begin{proof}[Proof of (b)]
    Suppose $a \in A$ is a unit. Then there 
    exists $b \in A$ so that 
    \[
        ab = 1_A. 
    \]
    Then $a,b \in A[x]$ as constant polynomials. 
    Once again 
    \[
        1_{A[x]} = 1_A = 
        ab = \phi(a)\phi(b) 
        = \phi(ab). 
    \]
\end{proof}
\begin{proof}
    From the above we have that 
    units in $A$ as constant polynomials 
    are units in $A[x]$.  

    Now let $f \in A[x]$ be a unit. Then 
    there is $g \in A[x]$ so 
    \[
        fg = 1. 
    \]
    Then since $\deg(g\cdot h) = \deg(g) + \deg(h)$ we have 
    \[
        \deg(f) + \deg(g) = \deg(f\cdot g) = \deg(1) = 0. 
    \]
    Then $\deg(f) = \deg(g) = 0$ so they are constant polynomials. 

    We require $A$ to be an integral domain because otherwise 
    we could take $A = \mathbb Z_4$ and $f = 2x + 1$. 
    Then $f^2 = 1$. 
\end{proof}

\begin{prob}
    Prove that a ring $A$ is a field if and only if 
    the following holds: $A \neq 0$ and every homomorphism 
    from $A$ to a non-zero ring is injective. 
\end{prob}
\begin{proof}
    Let $A$ and $B \neq 0$ be rings. 

    First assume $A$ is a field. 
    Then $\ker \phi$ is $A$ or $(0)$ for every homomorphism 
    $\phi : A \to B$. If $\ker \phi = A$ then 
    $\psi(x) = 0_B$ for all $x \in A$ and in particular 
    $1_B = \phi(1_A) = 0_B$ so $B$ is the trivial ring, but 
    this contradicts our assumption. Thus 
    $\ker \phi = (0)$, meaning $\phi$ is injective. Furtheremore, 
    by definition, $A \neq 0$.
    
    Now assume $A \neq 0$ and every homomorphism 
    $\phi : A \to B$ has $\ker \phi = \{0_A\}$ when $B$ is nonzero. 

    Let $I \subseteq A$ is some ideal. Consider the canonical 
    quotient homomorphism 
    \[
        \pi : A \to A/I \ \text{defined by} \ 
        \pi(a) = a+I. 
    \]
    Now either $A/I$ is nonzero or not. If it is the zero ring 
    then $\ker \pi = A$. 
    Otherwise if $A/I$ is nonzero, then by assumption 
    $\ker \pi = (0)$. By construction, 
    \[
        \ker \pi = I
    \]
    so every ideal $I$ of $A$ is either $(0)$ or $A$. Hence 
    $A$ is a field. 
\end{proof}

\begin{prob}
    Let $A$ be a ring, and let $x \in A$. Consider 
    the map $m_x : A \to A$ given by $m_x(y) = xy$ for 
    some $x \in A$. 
    \begin{enumerate}[label=(\alph*)]
        \item Show that $m_x$ is a ring homomorphism if 
            and only if $x = 1$. 
        \item Show that $x$ is a $0$-divisor if and only if 
            $m_x$ is non-injective. 
        \item Show that $x$ is a unit if and only if $m_x$ 
            is surjective. 
    \end{enumerate}
\end{prob}
\begin{proof}[Proof of (a)]
    Suppose $x = 1$. Then $m_x : A \to A$ is the 
    identity map $\operatorname{id} : A \to A$, 
    \[
        \operatorname{id}(y) = y,  
    \]
    which is a homomorphism. 

    Now suppose $m_x$ is a ring homomorphism and suppose 
    for a contradiction that $x \neq 1$. 
    Then 
    \[
        m_x(1) = x \neq 1 
    \]
    so $m_x$ does not preserve the multiplicative 
    identity and thus fails to be a homomorphism. 
\end{proof}

\begin{proof}[Proof of (b)]
    Suppose $x$ is a $0$-divisor, meaning $x \neq 0$ and there  
    exists $\xi \in A\setminus \{0\}$ so that 
    \[
        x \cdot \xi = 0.
    \]
    Then $\xi \in \ker m_x$ and $\xi \neq 0$ so 
    $m_x$ cannot be injective. 

    Suppose $m_x$ is not injective. 
    So $|\ker m_x| > 1$. Then there exists 
    $\xi \in \ker m_x \setminus \{0\}$. 
    Then 
    \begin{align*}
        m_x(\xi) &= m_x(0) \\ 
        x \cdot \xi = x \cdot 0 = 0
    \end{align*}
    so $x$ is a $0$-divisor.
\end{proof}
\begin{proof}[Proof of (c)]
    Suppose $x$ is a unit. Then there exists 
    $x^{-1} \in A$ so that 
    $xx^{-1} = 1$. 
    Then, for any $a \in A$ we see that 
    \[
        xx^{-1}a \overset{m_x}{\mapsto} a
    \]
    so $m_x$ is surjective. 

    Now suppose $m_x$ is surjective. Then, for 
    any $a \in A$ there exists $b \in A$ so that 
    \[
        m_x(b) = xb = a. 
    \]
    In particular this holds 
    for $1 \in A$ so there exists $x^{-1} \in A$ so that 
    \[
        m_x(x^{-1}) = xx^{-1} = 1. 
    \]
\end{proof}

\begin{prob}
    Let $A$ be a ring, and let $\{P_i\}_{i\in I}$ be a 
    non-empty chain of prime ideals (where "chain" ,eans 
    that for every two $i,j \in I$ either $P_i \subseteq P_j$ 
    or $P_j \subseteq P_i$). Prove that 
    $\bigcup_{i\in I}P_i$ and $\bigcap_{i\in I}P_i$
    are prime ideals. 
\end{prob}
\begin{proof}
    Let 
    \[
        J = \bigcup_{i \in I} P_i
    \]
    and 
    \[
        K = \bigcap_{i \in I} P_i.
    \]

    We begin with $J$. 
    Since the chain is nonempty, pick any $P_i$. Since 
    $P_i$ is an ideal, $0 \in P_i$ so $0 \in J \neq \emptyset$. 
    Let $x \in J$ and $r \in A$. Since $x \in J$, 
    $x \in P_i$ for some specific index $i$. Since $P_i$ is 
    an ideal we have $rx \in P_i$ so $rx \in J$. 
    Let $x,y \in J$. Then $x \in P_i$ and $y \in P_j$ for 
    some specific indicies $i,j\in I$. The chain-condition 
    gives us that $P_i \subseteq P_j$ or vice-versa. Assume 
    wlog the former holds. Then 
    $x,y \in P_j$ so $x+y \in P_j \subseteq J$. Thus 
    $J$ is closed under addition. 

    Now since every $P_i$ is prime we have $1 \not\in P_i$ for 
    every $i \in I$ so $1 \not \in J$. Thus $J \neq A$ and is 
    therefore a proper ideal. Let $a,b \in A$ such 
    that $ab \in J$. Then there is some $P_i$ so that 
    $ab \in P_i$. Since $P_i$ is prime we have 
    $a \in P_i$ or $b \in P_i$. 

    Now we do $K$. 
    Similar as before $0 \in P_i$ for all $i \in I$ so 
    $0 \in K \neq \emptyset$. 
    Let $x \in J$ and $r \in A$. Since $x \in J$ we have that 
    $x \in P_i$ for every $i$ so $ax \in P_i$ for every 
    $i$ hence $ax \in K$. Let $x,y \in K$. Then 
    $x,y \in P_i$ for every $i$ hence 
    $x +y \in P_i$ for every $i$ so $x+ y \in K$. 
    Thus $K$ is closed under addition. 

    Now supose $ab \in K$. Then $ab \in P_i$ for every 
    $i$ so $a \in P_i$ or $b \in P_i$ for every $i$. 
    Thus $a \in K$ or $b \in K$. 
\end{proof}

\begin{prob}
    Let $X$ be a topological space and $x \in X$. 
    We consider pairs $(U,f)$, where $U \subseteq X$ is an open 
    set with $x \in U$ and $f : U \to \mathbb R$ is a 
    continuous function. 

    Two pairs $(U,f)$ and $(V,g)$ are called equivalent if 
    there exists an open neighborhood $W$ of $x$ with 
    $W \subseteq U \cap V$ such that 
    $f\bigm|_W = g\bigm|_W$. An equivalence class is called 
    a germ of a function at $x$. We denote by $\mathcal O_{X,x}$ 
    the set of all such germs. 

    \begin{enumerate}[label=(\alph*)]
        \item Show that addition and multiplication of germs 
            are well defined by 
            \[
                (U,f) + (V,g) 
                := (U \cap V, f\bigm|_{U\cap V} 
                + g\bigm|_{U\cap V}), 
            \]
            \[
                (U,f) + (V,g) 
                := (U \cap V, f\bigm|_{U\cap V} 
                g\bigm|_{U\cap V}). 
            \]
            Verify that $\mathcal O_{X,x}$ is a ring 
            with these operations, and describe the 
            elements $0,1 \in \mathcal O_{X,x}$ as germs. 
        \item Show that a germ $(U,f)$ is a unit in 
            $\mathcal O(X,x)$ if and only if $f(x) \neq 0$. 
        \item Show that $\mathcal O_{X,x}$ is a local ring, 
            with maximal ideal 
            \[
                \mathfrak m = \{(U,f)\in \mathcal O_{X,x} 
                \bigm| f(x) = 0\}, 
            \]
            and residue field 
            \[
                \mathcal O_{X,x}/\mathfrak m \cong \mathbb R.
            \]
    \end{enumerate}
\end{prob}
\begin{proof}[Proof of (a)]
    Suppose $(U_1,f_1) = (U_2, f_2)$ and 
    $(V_1, g_1) = (V_2, g_2)$. By definition of the equivalence 
    relation 
    \begin{itemize}
        \item there exists an open neighborhood 
            $W_1 \subseteq U_1\cap U_2$ of $x$ such that 
            $f_1\bigm|_{W_1} = f_2\bigm|_{W_1}$, and 
        \item there exists an open neighborhood 
            $W_2 \subseteq V_1\cap V_2$ of $x$ such that 
            $g_1\bigm|_{W_2} = g_2\bigm|_{W_2}$.
    \end{itemize}
    Define $W = W_1 \cap W_2$. Since open neighborhoods 
    are closed under finite intersections we have that 
    $W$ is an open neighborhood of $x$. Furthermore, 
    $W \subseteq (U_1 \cap v_1) \cap (U_2 \cap V_2)$. 
    Restricting the sum to $W$ we get 
    \[
        (f_1 + g_2) \bigm|_W 
        = f_1\bigm|_W + g_1\bigm|_W 
        + f_2\bigm|_W + g_2\bigm|_W 
        = 
        (f_2 + g_2) \bigm|_W. 
    \]
    Thus addition is well-defined and this approach works  
    for multiplication as well. 

    The ring properties of associativity, commutativity and 
    distributivity are inherited from the pointwise operations 
    of $\mathbb R$. 

    The additive itentity is $(X,\textbf{0})$ where 
    $\textbf{0} : X \to \mathbb R$ is the constant 
    function $0$. 

    The multiplicative identity is 
    $(X, \textbf{1})$ where $\textbf{1} : X \to \mathbb R$ is 
    the constant $1$. 
\end{proof}
\begin{proof}[Proof of (b)]
    Suppose $(U,f)$ is a unit. Then there exists 
    $(V,g)$ such that $(U,f)\cdot (V,g) = (X,\textbf{1})$. 
    Then there is an open neighborhood $W \subseteq U \cap V$ 
    of $x$ where $f(k)g(k) = 1$ for all $k \in W$. 
    Evaluated at $x$ we get 
    \[
        f(x)g(x) = 1. 
    \]
    Since $f(x),g(x) \in \mathcal R$ we cannot have 
    $f(x) = 0$ as $0^{-1}$ does not exist.  

    Suppose $f(x) \neq 0$. Since $f: U \to \mathbb R$ is 
    continuous and $\{0\}^c = \mathbb R \setminus \{0\}$ is 
    an open subset of $\mathbb R$, the preimage 
    $V = f^{-1}(\{0\}^c)$ is open in $U$ and thus open 
    in $X$. 

    Since $f(x) \neq 0$ we have $x \in V$ so we can define 
    a continuous $g : V \to \mathbb R$ by 
    \[
        g(k) = \frac{1}{f(k)} \ \text{for all} \ k \in V.
    \]
    Then, on the open neighborhood $V$, we have 
    $f\bigm|_V \cdot g = 1\bigm|_V$. Thus 
    $[U,f]$ is a unit. 
\end{proof}

\newpage 

\section{A step back}
Everything thus far has been loosely covering 
Jørgen V. Rennemo's notes as I was planing 
to take this course at UiO as I usually do. 
Now I will be taking the equivalent 
MAT224 course at UiB instead! 
Thus I think it is fitting to go back, as 
there are some pretty big differences 
in the flavour of this course when compared with 
MAT4200 of UiO. 

Gathmann motivates commutative algebra 
by telling us that commutative rings 
provide a common algebraic language for things 
that initially look number-theoretic. In particular, 
polynomial equations define geometric objects, and 
the functions on those objects form commutative rings. 

The philosophy is that instead of studying geometric 
objects directly, study the ring of polynomial functions 
on it. This is quite a bit more algebraic geometry esque. 
Very cool! 

\begin{defn}
    Let $K$ be a field. The affine $n$-space over $K$ is 
    the set 
    \[
        \mathbb A^n_K = \{(a_1,\ldots,a_n) \bigm| a_i \in K\}.
    \]
    
\end{defn}
As a set, 
\[
    \mathbb A_K^n = K^n.
\]
The different notation emphasizes that we are not interested 
in the set as a vector space. We are thinking 
of it simply as a set of points on which 
polynomial equations can be imposed. 

For instance $\mathbb A^2_{\mathbb R}$
is simply the real plane. 

Let $f\in K[x_1,\ldots,x_n]$. 
Write 
\[
    f = \sum_{(i_1,\ldots,i_n) \in \mathbb N^n} 
    a_{i_1\ldots,i_n}x_1^{i_1}\cdots x_n^{i_n},
\]
where only finitely many coefficients are nonzero. 

Given $a = (a_1,\ldots,a_n) \in \mathbb A^n_K$, 
we evaluate $f$ at $a$ by 
\[
    f(a) = \sum a_{i_1, \ldots, i_n} a_1^{i_1} \cdots a_n^{i_n}.
\] 
The distinction is simply between the formal 
variables $x_j$ and a point with coordinates $a_j$ at 
which we substitute values.

\begin{defn}
    Let 
    \[
        S \subseteq K[x_1,\ldots,x_n]
    \]
    be any set of polynomials. Its zero locus is 
    \[
        V(S) = \{a \in \mathbb A^n_K \bigm| f(a) = 0 \ 
        \text{for every} \ f \in S\}.
    \]
    A subset of $\mathbb A^n_K$ that can be 
    written as $V(S)$ is called an \textbf{affine variety} 
    using Gathmann's terminology. 
\end{defn}

If $S = (f_1, \ldots, f_r)$ we normally write 
$V(f_1,\ldots,f_r)$. 

\begin{ex}
    Over $\mathbb R$, 
    \[
        V(x^2 + y^2 - 1) \subseteq \mathbb A^2_{\mathbb R}
    \]
    is the unit circle. 
    Indeed, 
    \[
        (a,b) \in V(x^2 + y^2 - 1)
    \]
    if and only if 
    \[ 
        a^2 + b^2 = 1. 
    \]
\end{ex}
\begin{ex}
    Consider 
    $$ 
        V(xy) \subseteq \mathbb A^2_{\mathbb K}.
    $$ 
    A point $(a,b)$ lies in $V(xy)$ iff 
    $ab = 0$. 
    Since $K$ is a field, this means 
    $a = 0$ or $b = 0$. 
    Hence 
    \[
        V(xy) = V(x) \cup V(y). 
    \]
    Geometrically this is the union of the $x$ and $y$ 
    axes. 
    This notion will appear later algebraically through 
    prime ideals and primary decomposition. How exciting!
\end{ex}

Now, why can't we simply go through equations themselves? 
Suppose $X = V(f)$. It would be tempting to say that 
$f$ is the equations of $X$. However, there is an immediate 
problem: 
\[
    V(f) = V(f^2). 
\]
Many entirely different collectioons of polynomials 
can have exactly the same zero locus. Gathmann uses 
this ambiguity to motivate replacing a 
chosen set of equations by all equations 
that vanish on the object. 

\begin{defn}
    Let $X \subseteq \mathbb A_K^n$. 
    We do not actuaklly need to assume $X$ is a variety 
    yet. 

    The ideal of $X$ is 
    \[
        I(X) = \{f \in K[x_1,\ldots,x_n] : f(a) = 0 \ 
        \text{for every} \ a \in X\}. 
    \]
\end{defn}

\begin{prop}
    \[
        I(x) \triangleleft K[x_1,\ldots,x_n]. 
    \]
\end{prop}
\begin{proof}
    Ceirtainly $0 \in I(X)$. 

    Let $f,g \in I(X)$. 
    For every $a\in X$, 
    \[ 
        (f-g)(a) = f(a) - g(a) = 0 - 0 = 0. 
    \]
    Hence $f - g \in I(X)$. 
    Thus $I(X)$ is an additive subgroup. 
    Now let $h \in K[x_1,\ldots,x_n]$ and 
    $f \in I(X)$. For every $a \in X$, 
    \[
        (hf)(a) = h(a)f(a) = h(a) \cdot 0 = 0. 
    \]
    Thus $hf \in I(X)$. 
    Therefore $I(X)$ is an ideal.
\end{proof}

We now take the quotient 
\[
    A(X) = K[x_1,\ldots,x_n]/I(X). 
\]
This is called the \textbf{coordinate ring} of $X$, or 
the ring of polynomials functions on $X$. 

Why quotient by $I(X)$? 

Suppose $f,g \in K[x_1,\ldots,x_n]$. 
In the quotient, $f + I(X) = g + I(X)$ if and only if 
$f-g \in I(X)$ so $f(a) = g(a)$ for every $a\in X$. 

Thus $f=g$ in $A(X)$ if and only if $f$ and $g$ define the 
same function on $X$. 

\begin{ex}
    Let $X = V(y-x^2) \subseteq \mathbb A^n_K$. 
    
    On $X$ we have $y = x^2$. Thus as polynomial functions on 
    $X$, y and $x^2$ are the same function. 
    
    Algebraically this is encoded by $A(X) = K[x,y]/I(X)$. 
\end{ex}  

There is some nuance here when working over finite fields, 
but Gathmann says the course will largely 
surpress this as we will be working over infinite fields 
like $\mathbb R$ and $\mathbb C$ for the most part. 

We will now introduce the ideal of a point. 

Take $a = (a_1,\ldots,a_n)$ in the affine $n$-space over some 
field $K$. Then 
\[
    I(a) = I(\{a\}) 
    = (x_1 - a_1, \ldots, x_n - a_n). 
\]

\begin{prop}
    For $a$ as above, we have 
    \[
        I(a) = (x_1 - a_1, \ldots, x_n - a_n). 
    \]
\end{prop}
\begin{proof}
    First suppose $f \in (x_1 - a_1, \ldots, x_n - a_n)$. 

    Then there exist polynomials $g_1,\ldots,g_n$ such that 
    \[
        f = \sum_{i=1}^n (x_i - a_i)g_i. 
    \]
    Evaluating at a, 
    \[
        f(a) = \sum_{i=1}^n (a_i - a_i)g_i(a) = 0. 
    \]

    Hence $f \in I(a)$. 
    Therefore 
    \[
        (x_1 - a_1,\ldots,x_n - a_n) \subseteq I(a). 
    \]

    For the reverse, let $f \in I(a)$, so 
    \[
        f(a_1,\ldots,a_n) = 0. 
    \]
    We can write 
    \begin{align*}
        f(x_1,\ldots,x_n) - f(a_1,\ldots,a_n)
        &= (f(x_1, x_2,\ldots,x_n) - f(a_1,x_2,\ldots,x_n)) \\ 
        &+ (f(a_1,x_2, \ldots, x_n) - f(a_1, a_2, x_3, \ldots, x_n)) \\ 
        &+ \cdots \\
        &+ (f(a_1,\ldots,a_{n-1},x_n) - f(a)). 
    \end{align*}
    The first summand is divisible by $x_1 - a_1$, the second 
    by $x_2 - a_2$, and so on. 
    Therefore 
    \[ 
        f - f(a) \in (x_1 - a_1, \ldots, x_n - a_n). 
    \]
    But $f(a) = 0$, so 
    \[ 
        f \in (x_1-a_1, \ldots, x_n - a_n)
    \]
    giving us the inclusion. 
    Together 
    \[
        I(a) = (x_1-a_1,\ldots,x_n - a_n).
    \]

\end{proof}

Immediately, the evaluation map 
\[ 
    \operatorname{ev}_a : K[x_1,\ldots,x_n] \to K, \ 
    f \mapsto f(a) 
\]
has 
\[
    \ker \operatorname{ev}_a = I(a). 
\]
It is surjective so 
\[
    \frac{K[x_1,\ldots,x_n]}{I(a)} \cong K. 
\]

So far our ideals lived in $K[x_1,\ldots,x_n]$. 
Suppose $X$ is already a variety. 

If $Y \subseteq X$, we can define 
\[
    I_X(Y) = \{f \in A(X) \bigm| f(y) = 0 \ \text{for every} \ 
    y \in Y\}. 
\] 
Conversely, if $S \subseteq A(X)$, define 
\[ 
    V_X(S) = \{x \in X \bigm| f(x) = 0 \ \text{for every} \ 
    f \in S\}. 
\]
We surpress subscripts when the variety is clear. 

\begin{prop}
    If $Y \subseteq Y' \subseteq X$, 
    then 
    \[ 
        I(Y') \subseteq I(Y). 
    \]
\end{prop}
The proof of this is quite short so I won't bother writing it. +

Likewise, if 
\[
    S \subseteq S' \subseteq A(X), 
\]
then 
\[
    V(S') \subseteq V(S). 
\]

Intuitively, larger geometry has a direct correspondence 
with a smaller ideal, and at the extremes 
\[ 
    \mathbb A^n_K = (0)
\]
in the usual infinite-field situation, while 
\[ 
    I(\emptyset) = K[x_1,\ldots,x_n] = (1),
\]
since every polynomials vacuously vanishes on the 
empty set. 

Suppose $Y \subseteq X$. 

We can form 
\[
    Y \mapsto I(Y) \mapsto V(I(Y)). 
\]

Ceirtainly, 
\[
    Y \subseteq V(I(Y)).
\]
Every polynomials in $I(Y)$ was defined precisely 
to vanish on $Y$. 
Similarly, for $S\subseteq A(X)$, 
\[
    S \subseteq I(V(S)). 
\]

If $Y$ is already a subvariety, something stronger holds. 

\begin{prop}
    If $Y \subseteq X$ is a subvariety, then 
    \[
        V(I(Y)) = Y. 
    \]
\end{prop}
\begin{proof}
    We know $Y \subseteq V(I(Y))$. 
    Because $Y$ is a subvariety, there exists some 
    $S \subseteq A(X)$ such that 
    \[
        Y = V(S). 
    \]
    Every element of $S$ vanishes on $Y$, so 
    \[
        S \subseteq I(Y). 
    \]
    Since taking zero loci reverses inclusion, 
    \[ 
        V(I(Y)) \subseteq V(S) = Y. 
    \]
\end{proof}

Now for a really nice result. 

Let $Y \subseteq X$ be a subvariety. 

Any polynomial function on $X$ can be restricted to $Y$. 
Thus we get a ring homomorphism 
\[
    \rho A(X) \to A(Y), \ f \mapsto f\bigm|_{Y}. 
\]

\begin{prop}
    The map $\rho$ is surjective and 
    \[
        \ker \rho = I(Y). 
    \]
    Hence 
    \[
        A(X) / I(X) \cong A(Y).
    \]
\end{prop}
\begin{proof}
    By the definition of polynomial functions on 
    $Y$, every polynomial function on $Y$ is obtained by 
    restricting a polynomial function from the ambient 
    space, so $\rho$ is surjective. 
    Furthermore, 
    \[
        f \in \ker \rho 
    \]
    if and only if 
    \[ 
        f\bigm|_Y = 0, 
    \]
    which means precisely that 
    \[
        f(y) = 0
    \]
    for every $y\in Y$. 

    Thus $\ker \rho = I(Y)$. 
    The first isomorphism theorem then gives 
    \[
        A(X)/I(Y) 
        = \frac{K[x_1,\ldots,x_n]/I(X)}{I(Y)} 
        \cong A(Y).
    \]
\end{proof}

The last major thing now concerns maps between 
varieties. 

Let $X \subseteq \mathbb A_K^n$ and $Y \subseteq \mathbb A_K^m$
be varieties. 

\begin{defn}
    A \textbf{morphism of varieties} 
    \[
        f : X \to Y
    \]
    is a map given by polynomials: 
    \[
        f(x) = (f_1(x), \ldots, f_m(x)),
    \]
    where 
    \[
        f_i \in K[x_1,\ldots,x_n]. 
    \]
    Gathmann then associates to $f$ a ring homomorphism 
    \[
        f^* : A(Y) \to A(X)
    \]
    defined by 
    \[
        f^*(g) = g \circ f. 
    \]
    Notice the direction: 
    \[
        X \overset{f}{\to} Y 
    \]
    produces 
    \[
        A(Y) \overset{f^*}{\to} A(X). 
    \]
\end{defn}

The cool thing now is that the arrow reverses. 

Suppose you have an ordinary function 
$f: X \to Y$ and a real-valued function $g : Y \to \mathbb R$. 
Composition gives $g \circ f : X \to \mathbb R$. 

So a map of spaces $X \to Y$ naturally lets you 
pull functions backwards 
\[
    \{\text{functions on} \ Y\} 
    \to 
    \{\text{functions on} \ X\}.
\]
This is called \textbf{pullback}. 

Hence the notation $f^*$. 
This contravariance is very important further into algebraic 
geometry. 

\subsection{A few practice problems}

\begin{prob}
    Determine geometrically 
    \[
        V(xy, xz) \subseteq \mathbb A_K^3. 
    \]
    Try to describe it as a union of familiar linear pieces. 
\end{prob}

\begin{proof}[Solution]
    We have 
    \[
        V(xy, xz) = \{(x,y,z) \in K^3 : xy = 0, xz = 0\}. 
    \]
    Both equations share a common factor of $x$. 
    If $x = 0$, then both equations vanish, so 
    this gives the plane 
    \[
        V(x). 
    \]
    If $x \neq 0$, then since $K$ is a field we may cancel 
    $x$ to get $y = 0$, $z = 0$, giving us the 
    $x$-axis $V(y,z)$. Therefore 
    \[
        V(xy,xz) = V(x) \cup V(y,z).
    \]
    Geometrically this is the union of the 
    $yz$-plane and $x$-axis.
\end{proof}

\begin{prob}
    Let $a = (2,-1) \in \mathbb A^2_{\mathbb R}$.
    Without a long proof, determine $I(a)$ and identify 
    \[
        \mathbb R[x,y]/I(a)
    \]
\end{prob}
\begin{proof}[Solution]
    We showed earlier that 
    $$
        I(a) = (x-2, y + 1).
    $$
    Define the ring homomorphism 
    $$
        \phi : \mathbb R[x,y] \to \mathbb R, \
        f \mapsto f(2,-1). 
    $$
    The map is surjective and $\ker \phi = I(a)$. 
    Thus 
    \[
        \mathbb R[x,y] / I(a) \cong \mathbb R.
    \]
\end{proof}

\section{Operations on ideals}

Let $R$ be a commutative ring with $1 \in R$. 

\begin{defn}
    Let $f_1,\ldots,f_n \in R$. The ideal generated by 
    $f_1,\ldots,f_n$ is 
    \[
        (f_1,\ldots,f_n) = \left\{\sum_{i=1}^n a_if_i \biggm| 
        a_1,\ldots,a_n \in R\right\}. 
    \]
    It is the smallest ideal containing all of 
    \[
        f_1,\ldots,f_n. 
    \]
\end{defn}
The coefficients are allowed to be arbitrary elements of the 
ring. For instance, in $K[x,y]$, 
\[
    (x,y) = \{fx+gx \bigm| f,g \in K[x,y]\},
\]
meaning 
\[
    x^4 + x^2y + y^7 \in (x,y). 
\]

For principal ideals (ideals generated by one elements), 
\[
    (a) = \{ra \bigm| r \in R\}. 
\]
If $a \bigm| b$, then $b = ra$ for some $r \in R$ so 
\[
    b \in (a). 
\]
In fact, 
\[
    a \bigm| b \Rightarrow (b) \subseteq (a).
\]

\subsection{Redundant generators}

Suppose $I = (f_1,\ldots,f_n)$. If one generator already 
belongs to the ideal generated by the others, it can (and should) 
be ommited. 

For example, 
\[
    (x,y,xy) = (x,y)
\]
because 
\[
    xy = y \cdot x \in (x).
\]

\subsection{Sum of ideals}
Let $I, J \triangleleft R$. 

\begin{defn}
    \[
        I + J 
        = 
        \{a+b \bigm| a \in I, \ b \in J\}.
    \]
\end{defn}
$I+J$ is the smallest ideal such that $I, J$ are contained within it.

Generally, if 
\[
    I = (f_1,\ldots,f_n)
\]
and 
\[
    J = (g_1,\ldots,g_n), 
\]
then 
\[
    I + J = (f_1,\ldots,f_n, g_1,\ldots,g_n).
\]
It is prudent to remove redundant generators after concatenating. 

For instance, 
\[
    (x^2, xy) + (y) = (x^2, xy, y)
\]
but 
\[
    xy \in (y)
\]
so we simplify to 
\[
    (x^2, y).
\]

\subsection{Intersection of ideals}
The intersection 
\[
    I \cap J
\]
is the set of elements which belong to both $I,J$, meaning 
it's just the usual intersection of sets operation. 
Unlike sums, intersections are often harder to compute directly 
from generators. 

Let $R = K[x,y]$. Suppose $f \in (x) \cap (y)$. Then 
\[
    f = ax, \ a \in R
\]
and 
\[
    f = by, \ b \in R.
\]
Since $x,y$ are relatively prime in the UFD $K[x,y]$ we have 
that $xy \bigm| f$. Therefore $f \in (xy)$ giving us that 
\[
    (x) \cap (y) \subseteq (xy). 
\]
Conversely, $f = hxy = (hx)y = (hy)x$, so $f \in (x)$ and $f \in (y)$.
Thus 
\[
    (xy) \subseteq (x) \cap (y).
\]

For monomials $m_1,m_2$ we have 
\[
    (m_1) \cap (m_2) = (\operatorname{lcm}(m_1,m_2)). 
\]

\subsection{Product of ideals}

\begin{defn}
    For ideals $I,J \triangleleft R$, 
    \[
        IJ
    \]
    is the ideal generated by all products $ab$ with $a \in I$ 
    and $b \in J$. 
    Equivalently, 
    \[
        IJ
        = 
        \left\{
            \sum_{k=1}^n a_kb_k \biggm| a_k \in I, \ b_k \in J
        \right\}.
    \]  
\end{defn}
It is important that the sums are finite, 
\[
    \{ab : a \in I, \ b \in J\}
\]
need not be an ideal. 

A nice rule is that if $I = (f_1,\ldots,f_r)$ and 
$J = (g_1,\ldots,g_s)$ then 
\[
    IJ = (f_ig_j : i \in [r], \ j \in [s]). 
\]

For instance 
\[
    (x,y)^2 = (x,y)(x,y) = (x^2, xy, y^2)
\]
remembering that $xy = yx$. 

\begin{prop}
    For ideals $I,J \triangleleft R$, 
    \[
        IJ \subseteq I \cap J. 
    \]
\end{prop}
\begin{proof}
    Let $a \in I$ and $b \in J$. Then, since $I,J$ are ideals 
    we have $ab = ba \in I$ and $ab \in J$. Thus 
    $ab \in I \cap J$. Therefore every generator belongs to 
    $I \cap J$ and 
    \[
        IJ \subseteq I \cap J.
    \]
\end{proof}

For powers of ideals we define it as follows. 
We define $I^2 = II$, $I^3 = III$ and generally 
\[
    I^n = \prod_{i\in[n]} I = 
    \underbrace{I\cdots I}_{n \ \text{factors}}.
\]

\subsection{Ideal quotient}

\begin{defn}
    Let $I,J \triangleleft R$. The ideal quotient is 
    \[
        I : J = \{r \in R \bigm| rJ \subseteq I\}. 
    \]
\end{defn}
For principal ideal $J = (g)$ we have that 
\[
    rJ \subseteq I
\]
precisely when 
\[
    rg \in I. 
\]

For example if we take $(xy) : (x)$ we want all 
$r \in R$ such that 
\[
    rx \in (xy),
\]
meaning 
\[
    rx = hxy
\]
for some $h \in K[x,y]$. 
Since $K[x,y]$ is a domain we can cancel $x$ to get 
\[
    r = hy. 
\]
Thus $r \in (y)$. Conversely if 
\[
    r = hy, 
\]
then $rx = hxy \in (xy)$. 
So, maybe unsurprisingly, 
\[
    (xy) : (x) = (y).
\]

\subsection{Quotient by nonprincipal ideal}
If $J = (g_1,\ldots,g_s)$, then 
\[
    rJ \subseteq I
\]
means 
\[
    rg_i \in I
\]
for every generator $g_i$. 
Thus 
\[
    I : (g_1,\ldots,g_s) = \bigcap_{i=1}^s (I:(g_i)). 
\]

\subsection{Radical of an ideal}
\begin{defn}
    For an ideal $I \triangleleft R$, 
    \[
        \sqrt I = \{r \in R \bigm| r^n \in I \ \text{for some} \ 
        n \geq 1\}.
    \]
\end{defn}
The radical consists of elements which lie in $I$ for sufficiently 
large power. Clearly 
\[
    I \subseteq \sqrt I.
\]

\begin{defn}
    An ideal is radical if 
    \[
        \sqrt I = I. 
    \]
    Equivalently, 
    \[
        r^n \in I \Rightarrow r \in I. 
    \]
\end{defn}

A non-example is $(x^2)$ since $x \not\in (x^2)$. 
Interestingly 
\[
    \sqrt{(x^{100})} = \sqrt{(x^2)} = (x).
\]

\subsection{Nilpotents}

\begin{defn}
    An element $a \in R$ is nilpotent if 
    \[
        a^n = 0
    \]
    for some $n \geq 1$. 
    The set of nilpotents is thus 
    \[
        \sqrt{(0)},
    \]
    usually called the nilradical of $R$.

    A ring $R$ is called reduced if it has no nonzero nilpotent 
    elements. Equivalently, 
    \[
        R \ \text{reduced} \Leftrightarrow \sqrt{(0)} = (0).
    \]
\end{defn}

\begin{prop}
    Let $I \triangleleft R$. Then 
    \[
        I \ \text{is radical} \Leftrightarrow 
        R/I \ \text{is reduced}. 
    \]
\end{prop}
\begin{proof}
    An element $\overline a = a + I \in R/I$ is nilpotent 
    if and only if $\overline a^n = 0$ for some $n \geq 1$.
    This is equivalent to $a^n \in I$. 
    Thus $R/I$ has no nonzero nilpotents exactly when 
    $a^n \in I \Rightarrow a \in I$ which is precisely the statement 
    that $I $ is radical. 
\end{proof}

There are a lot of these links between behaviours of the ideal 
and the quotient, two of which from a first course in algebra 
is the correspondence: 
\[
    I \ \text{prime} \Leftrightarrow R/I \ \text{domain}, 
\]
and 
\[
    I \ \text{maximal} \Leftrightarrow R/I \ \text{field}. 
\]

\subsection{Geometry of radicals}
For an ideal $I \triangleleft K[x_1,\ldots,x_n]$, we have 
\[
    V(I) = V(\sqrt{I}).
\]
\begin{proof}
    Since $I \subseteq \sqrt I$, taking zero loci reverses 
    the inclusion 
    \[
        V(\sqrt I) \subseteq V(I). 
    \]
    Conversely let $x \in V(I)$. Take any $f \in \sqrt{I}$. 
    Then $f^m \in I$ for some $m$. Therefore 
    $f^m(x) = f(x)^m = 0$. Since $K$ is a field we get 
    \[
        f(x) = 0.
    \]
    Therefore $x \in V(\sqrt I)$, meaning 
    \[
        V(I) \subseteq V(\sqrt I). 
    \]
    All in all: 
    \[
        V(I) = V(\sqrt I).
    \]
\end{proof}

Let $Y \subseteq \mathbb A_K^n$. Recall that 
\[
    I(Y) = \{f : f(y) = 0 \ \forall y \in Y\}. 
\]
Then: 
\begin{prop}
    The ideal $I(Y)$ is radical. 
\end{prop}
\begin{proof}
    Suppose $f^n \in I(Y)$. Then 
    \[
        f(y)^n = 0, \ \forall y \in Y. 
    \]
    Since $K$ is a field, 
    \[
        f(y) = 0. 
    \]
    Therefore $f \in I(Y)$, meaning $I(Y)$ is radical. 
\end{proof}

\subsection{Other stuff}

\begin{prop}
    \[
        \sqrt{IJ} = \sqrt{I\cap J} = \sqrt{I} \cap \sqrt J.
    \]
\end{prop}

\begin{prop}
    \[
        V(I + J) = V(I) \cap V(J).
    \]
\end{prop}

\begin{prop}
    \[
        V(IJ) = V(I \cap J) = V(I) \cup V(J).
    \]
\end{prop}

\begin{defn}
    Ideals $I,J$ are called coprime or comaximal if 
    \[
        I+J = R. 
    \]
\end{defn}

\begin{prop}
    If $I + J = R$, then 
    \[
        IJ = I \cap J.
    \]
\end{prop}


\begin{thm}[Chinese remainder theorem]
    Let $I_1,\ldots, I_n \triangleleft R$. 
    Define 
    \[
        \phi : R \to \prod_{i=1}^n R/I_i
    \]
    by 
    \[
        r \mapsto (r + I_1, \ldots, r + I_n). 
    \]
    Then 
    \[
        \ker \phi = \bigcap_{i=1}^n I_i.
    \]
    If $I_i$ are pairwise coprime, then $\phi$ is surjective. 
    Thus, by the first isomorphism theorem: 
    \[
        R\big/\bigcap_{i=1}^n I_i \cong \prod_{i=1}^n R/I_i. 
    \]

\end{thm}

\subsection{Problems}

Let $R = K[x,y]$.
\begin{prob}
    Compute 
    \[
        (x^3, y ) + (x, y^4). 
    \]
\end{prob}
\begin{proof}[Solution]
    \[
        (x^3,y) + (x,y^4) = (x^3, y, x, y^4). 
    \]
    Since $x^3 \in (x)$ and $y^4 \in y$ it simplifies to 
    $(x,y)$. We could also observe that $(x^n) \subseteq (x)$ 
    for any $n$ (tying it back to radicals).
\end{proof}

\begin{prob}
    Compute 
    \[
        (x^2y) \cap (xy^3). 
    \]
\end{prob}
\begin{proof}[Sol]
    Suppose $f \in (x^2y) \cap (xy^3)$. Then 
    \begin{align*}
        f = gx^2y &= hxy^3 \\
    \end{align*}
    Both are divide by $x^2y^3$ which is 
    $\operatorname{lcm}(x^2y, xy^3)$ in $R$.
    Thus $f \in (x^2y^3)$, yielding 
    \[
        (x^2y) \cap (xy^3) \subseteq (x^2y^3). 
    \]
    Now suppose $f \in (xy)$. Then 
    \begin{align*}
        f &= ax^2y^3 \\ 
          &= ax(xy^3) \\ 
          &= ay^2(x^2y). 
    \end{align*}
    Thus 
    \[
        (x^2y) \cap (xy^3) \supseteq (x^2y^3). 
    \]
\end{proof}

\begin{prob}
    Compute 
    \[
        (x,y)(x^2,y^2)
    \]
\end{prob}
\begin{proof}[Sol]
    Recall that 
    \[
        (x,y)(x^2,y^2) 
        = (xx^2, xy^2, yx^2, yy^2) 
        = (x^3, xy^2, yx^2, y^3). 
    \]
\end{proof}

\begin{prob}
    Compute 
    \[
        (x^3y^2) : (x^2). 
    \]
\end{prob}
\begin{proof}[Sol]
    We want to find all $r \in R$ so that 
    $rx^2 \in (x^3y^2)$. 
    Then 
    \begin{align*}
        rx^2 &= hx^3y^2 \\ 
        r &= hxy^2. 
    \end{align*}
    Therefore 
    \[
        (x^3y^2) : (x^2) = (xy^2). 
    \]
\end{proof}

\begin{prob}
    Compute 
    \[
        (x^2, xy) : (x). 
    \]
\end{prob}
\begin{proof}[Sol]
    We want to find $r \in R$ so that 
    \[
        rx = hx^2 
    \]
    and 
    \[
        rx = hxy 
    \]
    similtaneously. 
    Proceeding as earlier we get 
    \[
        (x) \cap (y) = (xy). 
    \]
\end{proof}

\begin{prob}
    Compute 
    \[
        \sqrt{(x^8, y^{10})}. 
    \]
\end{prob}
\begin{proof}[Sol]
    We want to find all $r \in R$ so that 
    \[
        r^n \in (x^8, y^10), \ \text{for some} \ n \geq 1. 
    \]
    This means 
    \[
        r^n = ax^8 + by^{10}.
    \]
    We claim $(x,y) \subseteq \sqrt{(x^8, y^{10})}$. 
    For $x$ choose $n = 8$ and similarly for $y$ choose 
    $n=10$. Thus $x,y \in \sqrt{(x^8, y^{10})}$. 
    This yields the inclusion: 
    \[
        (x,y) \subseteq \sqrt{(x^8, y^{10})}. 
    \]
    Let $f(x,y)$ be some polynomial in $\sqrt{(x^8, y^{10})}$. 
    Then there exists some $m$ so that 
    \[
    f^m \in (x^8, y^{10}). 
    \]
    We have 
    \[
        f^m = ax^8 + by^{10}
    \]
    Since $x,y \equiv 0$ modulo $(x,y)$ we have 
    \[
        f^m = 0
    \]
    and since $K$ is a field we then get 
    \[
        f(x,y) = 0
    \]
    so 
    \[
        f \in (x,y). 
    \]
\end{proof}

\begin{prob}
    Decide whether 
    \[
        (x) + (x-1) = K[x]. 
    \]
    If yes, use it to compute 
    \[
        (x) \cap (x-1).
    \]
\end{prob}
\begin{proof}
    Recall that 
    \[
        I = (x) + (x-1) = (x, x-1). 
    \]
    Then $x - (x - 1) = 1 \in I$. 
    Now for any $f \in K[x]$ we have 
    $f \cdot 1 = f \in I$ so $I = K[x]$. 

    Since $(x)$ and $(x-1)$ are comaximal we have 
    $(x) \cap (x-1) = (x)(x-1) = (x^2 - x)$ . 
\end{proof}

\begin{prob}[Set 2, problem 4 (Gathmann 1.8)]
    Let $I_1,\ldots,I_n$ be pairwise coprime ideals 
    in $R$. Prove that 
    \[
        I_1\cdots I_n = I_1 \cap \cdots \cap I_n. 
    \]
\end{prob}
\begin{proof}
    We have that for $i \neq j$, 
    \[
        I_i + I_j = R.
    \]

    $I_1 \cdots I_n \subseteq I_1 \cap \cdots \cap I_n$ is immediate. 
    Now, for $n = 2$, suppose $f \in I_1 \cap I_2$. 
    Then, since $I_1 + I_2 = R$, there exists $a \in I_1, \ b \in I_2$
    so that 
    \[
        a+b = 1. 
    \]
    Multiplying by $x$ we get 
    \[
        ax+bx = x. 
    \]
    Since $x \in I_2$ and $a \in I_1$ we have 
    $ax \in I_1I_2$. Likewise, since $x \in I_1$ and $b \in I_2$ 
    we have $bx \in I_1I_2$. 
    Thus $x = ax + bx \in I_1I_2$. Hence 
    \[
        I_1I_2 \supseteq I_1\cap I_2. 
    \]

    For the induction step, assume $n > 2$ and that the claim 
    holds for $n - 1$ pairwise coprime ideals. 

    First, we observe that if $I_1$ is coprime to each $I_k$ for 
    $k = 2, \ldots, n$, then $I_1 + (I_2 \cdots I_n) = R$. Indeed, 
    for each $k \ge 2$, choose $a_k \in I_1$ and $b_k \in I_k$ such 
    that $a_k + b_k = 1$. Expanding the product 
    \[
        1 = \prod_{k=2}^n (a_k + b_k)
    \]
    yields $1 = A + b_2 \cdots b_n$, where $A \in I_1$ 
    (since every term in the expansion except $b_2 \cdots b_n$ 
    contains at least one factor $a_k \in I_1$). 
    Since $b_2 \cdots b_n \in I_2 \cdots I_n$, this shows that 
    $I_1 + (I_2 \cdots I_n) = R$.

    Now, let $J = I_2 \cdots I_n$. By the induction hypothesis, 
    $J = I_2 \cap \cdots \cap I_n$. Since $I_1$ and $J$ are coprime, 
    applying the $n = 2$ base case to $I_1$ and $J$ yields
    \[
        I_1 \cdots I_n = I_1 J = I_1 \cap J = I_1 \cap (I_2 \cap 
        \cdots \cap I_n) = I_1 \cap \cdots \cap I_n,
    \]
    completing the induction.
\end{proof}

\subsection{Some problems from set 3}

\begin{prob}[Gathmann 1.19]
    \medskip 
    \  

    Let $\phi : R \to R'$ be a ring homomorphism. 

    We denote $I^c$ to be the inverse image of an 
    ideal or contraction of $I$. 

    For $I \triangleleft R$ the ideal generated 
    by the image of $\phi[I]$ is called the image 
    ideal or extension by $\phi$, written as 
    $I^e$ if the morphism is clear from context. 

    Prove: 

    (c) $(IJ)^e = I^eJ^e$ for all ideals $I,J$ of $R$. 
    
    (d) $(I\cap J)^c = I^c \cap J^c$ for all ideals 
    $I,J$ of $R'$. 
\end{prob}

\begin{proof}[Proof of (c)]
    Let $I,J$ be ideals of $R$. Let $\phi : R \to R'$ be a 
    homomorphism of ring. 

    We begin by proving the forward inclusion. 

    Suppose $x \in IJ$, meaning $\phi(x) \in (IJ)^e$. 
    By definition, 
    \[
        x = \sum_{i} \lambda_ix_i
    \]
    where $\lambda_i \in R, \ x_i \in IJ$. 
    Then, 
    \begin{align*}
        \phi\left(\sum_i \lambda_i x_i\right)
        &= \sum_i \phi(\lambda_i)\phi(x_i).
    \end{align*}
    For each $i$, we have that $\phi(x_i) \in IJ$
    and $\phi(\lambda_i) \in R'$. 
    Thus each $\lambda(x_i) \in I^eJ^e$. Since $I^eJ^e$ absorbs
    $\phi(\lambda_i)$ we get that each 
    $\phi(\lambda_i)\phi(x_i) = \phi(\lambda_ix_i) \in I^eJ^e$.
    As the ideal is closed under addition we get that 
    \[
        \phi(x) = \phi\left(\sum_i \lambda_i x_i\right)
        \in I^eJ^e
    \]
    and as $x$ was arbitrarily chosen we finally get 
    \[ 
        (IJ)^e \subseteq I^eJ^e. 
    \]

    Now for the reverse inclusion.

    Suppose $y \in I^eJ^e$. Then 
    \[
        y = \sum_k u_k v_k, 
    \]
    where $u_k \in I^e, \ v_k \in J^e$. 
    Now, because $I^e, J^e$ are generated by 
    $\phi[I], \ \phi[J]$, write 
    \[
        u_k = \sum_i r'_{ki}\phi(a_{ki}), \ 
        v_k = \sum_j s'_{ki}\phi(b_{ki}), 
    \]
    with $a_{ki} \in I, \ b_{ki} \in J$ and 
    $r'_{ki}, s'_{ki} \in R'$. 

    Then 
    \begin{align*}
        y &= \sum_{k} \left(\sum_i r'_{ki}\phi(a_{ki})\right)
        \left(\sum_j s'_{ki}\phi(b_{kj})\right) \\ 
        &= \sum_{k,i,j}r'_{ki}s'_{kj}\phi(a_{ki})\phi(b_{kj}) \\
        &= \sum_{k,i,j}r'_{ki}s'_{kj}\phi(a_{ki}b_{kj}). 
    \end{align*}
    But 
    \[
        a_{ki}b_{kj} \in IJ, 
    \]
    so each 
    \[
        \phi(a_{ki}b_{kj})
    \]
    is one of the generators of $(IJ)^e$. Hence 
    \[
        y \in (IJ)^e. 
    \]
    Therefore, 
    \[
        I^eJ^e \subseteq (IJ)^e. 
    \]
\end{proof}
\begin{proof}[Proof of (d)]
    Let $I,J\triangleleft R'$. For any $x\in R$, we have
    \begin{align*}
    x\in (I\cap J)^c
    &\iff \phi(x)\in I\cap J\\
    &\iff \phi(x)\in I \text{ and } \phi(x)\in J\\
    &\iff x\in I^c \text{ and } x\in J^c \\
    &\iff x\in I^c\cap J^c.
    \end{align*}
    Therefore,
    \(        (I\cap J)^c=I^c\cap J^c.
       \)
\end{proof}

\begin{prob}[Gathmann 1.22]
    Let $I \subset J$ be ideals in a ring $R$. 
    By Lemma 1.21 (in Gathmann's notes), the extension 
    $J/I$ of $J$ by the quotient map $R \mapsto R/I$
    is an ideal in $R/I$. Prove that 
    \[
        (R/I)/(J/I) \cong R/J. 
    \]
    
\end{prob}
\begin{proof}
    Let 
    \[
        \psi : R/I \to R/J
    \]
    be defined by 
    \[
        \psi(r + I) = r + J.
    \]
    We first check that $\psi$ is well-defined. 

    Suppose 
    \[
        r + I = s + I. 
    \]
    Then 
    \[  
        r-s \in I. 
    \]
    Since $I \subset J$, we also have 
    \[
        r-s \in J. 
    \]
    Hence 
    \[
        r + J = s + J. 
    \]
    Thus $\psi$ is well-defined. 
    It is clearly a ring homomorphism. 
    Moreover, $\psi$ is surjective since every element 
    of $R/J$ has the form  
    \[
        r+J = \psi(r + I)
    \]
    for some $r \in R$. 
    Computing the kernel we have 
    \[
        r + I \in \ker \psi
    \]
    if and only if 
    \[
        r + J = J. 
    \]
    Equivalently, 
    \[
        r \in J. 
    \]
    Therefore 
    \[
        \ker \psi = \{r + I \bigm| r \in J\} = J/I. 
    \]
    Hence, by the First Isomorphism Theorem, 
    \[
        (R/I)/(J/I) = (R/I)/\ker\psi \cong \im \psi. 
    \]
    As $\psi$ is surjective, 
    \[
        (R/I)/(J/I) \cong R/I. 
    \]
\end{proof}

\begin{prob}
    View $\mathbb Z$ as a $\mathbb Z$-module. Prove that 
    $2\mathbb Z$ is a submodule. 

    Then, determine whether $2\mathbb Z \cup 3\mathbb Z$ is 
    a submodule of $\mathbb Z$.
\end{prob}
\begin{proof}[Proof \& solution]
    We have that $2\mathbb Z = \{\ldots, -2, 0, 2, \ldots\} \neq 
    \emptyset$. 
    Furthermore take any $x,y \in 2 \mathbb Z$. Then 
    there are some $n,m$ so that 
    \[
        x = 2n, \ y = 2m
    \]
    so 
    \[
        x - y = 2n - 2m = 2(n-m) \in 2\mathbb Z
    \]
    hence the subgroup criterion is satisfied. 
    Now take $r \in \mathbb Z$. 
    Then 
    \[
        r \cdot x = r \cdot 2n = 2(rn) \in 2\mathbb Z. 
    \]
    Thus $2\mathbb Z$ is a $\mathbb Z$-module and, of course 
    \[
        2\mathbb Z \subseteq \mathbb Z.
    \]

    The latter is not a submodule as it is not even a group: 
    \[
        2 + 3 = 5 \not\in2\mathbb Z \cup 3\mathbb Z.
    \]
\end{proof}

\begin{prob}
    Let 
    \[
        M = \mathbb Z/12\mathbb Z
    \]
    viewed as a $\mathbb Z$-module, and let 
    \[
        N = (\overline 4). 
    \]
    \begin{enumerate}
        \item Determine $N$, 
        \item Determine the elements of the quotient module 
            \[
                M/N,
            \]
        \item Find a familiar $\mathbb Z$-module isomorphic 
            to $M/N$.
    \end{enumerate}
\end{prob}
\begin{proof}[Sol]
    \[
        (\overline 4) = \{\overline 0,\overline 4,\overline 8\}. 
    \]
    \[
        M/N = \{m + N \bigm| m = 0,1,2,\ldots,12\} 
        = \{\overline m + N \bigm| \overline m = 0,1,2,3\}. 
    \]

    Let $\phi : M \to N \to \mathbb Z_4$ by 
    $\phi(m) = m \mod 4$. Then $\ker \phi = N$ and it is 
    indeed surjective so the first isomorphism theorem gives us 
    \[
        M/N \cong \mathbb Z_4. 
    \]
\end{proof}

\begin{prob}
    View 
    \[
        M = \mathbb Z/12\mathbb Z
    \]
    as a $\mathbb Z$-module. 
    Find a composition series 
    \[
        0 = M_0 \subsetneq M_1 \subsetneq \cdots \subsetneq M_n = M.
    \]
    For each $i$ determine the quotient 
    \[
        \frac{M_i}{M_{i-1}}.
    \]
    Finally, use this to determine $\ell(M)$.
\end{prob}
\begin{proof}
    We find the series 
    \[
        0 \subsetneq \mathbb Z/6\mathbb Z \subsetneq 
        \mathbb Z/2\mathbb Z \subsetneq \mathbb Z/12\mathbb Z.
    \]
    We have 
    \[
        \frac{\mathbb Z/12\mathbb Z}{\mathbb Z/2\mathbb Z}
        \cong \mathbb Z/2\mathbb Z, 
    \]
    \[
        \frac{\mathbb Z/2\mathbb Z}{\mathbb Z/6\mathbb Z} 
    \]
    has size $6/2 = 3$ and since $\mathbb Z$-modules are abelian 
    groups we must have that it is isomorphic to $\mathbb Z_3$. 
    Lastly 
    \[
        \frac{\mathbb Z/6\mathbb Z}{(0)} \cong \mathbb Z/6\mathbb Z 
        \cong 2 \mathbb Z.
    \]
    All of the above are simple so this is a composition series 
    with $\ell(M) = 3$. 
\end{proof}

\begin{prob}
    Suppose $M$ has finite length and $N \subseteq M$ is a 
    submodule. Assume you already know that both $N$ and 
    $M/N$ have finite length. 
    Suppose 
    \[
        0 = N_0 \subsetneq N_1 \subsetneq \cdots \subsetneq N_r = N
    \]
    is a composition series for $N$, and 
    \[
        0 = Q_0 \subsetneq Q_1 \subsetneq \cdots \subsetneq Q_s = M/N
    \]
    is a composition series for $M/N$. 
    Let 
    \[
        \pi : M \to M/N
    \]
    be the quotient map. 

    Construct a composition series for $M$ using the $N_i's$ inverse 
    images 
    \[
        \pi^{-1}[Q_j]. 
    \]
    Deduce then that 
    \[
        \ell(M) = \ell(N) + \ell(M/N).
    \]
\end{prob}
\begin{proof}
    For each $j \in \{0, 1, \dots, s\}$, since $Q_j$ is a submodule of $M/N$, its preimage under the quotient map $\pi$, denoted $\pi^{-1}[Q_j]$, is a submodule of $M$ containing $\ker \pi = N$. 

    Evaluating the preimages at the endpoints of the composition series for $M/N$, we have:
    \[
        \pi^{-1}[Q_0] = \pi^{-1}[\{0\}] = \ker \pi = N, \quad \text{and} \quad \pi^{-1}[Q_s] = \pi^{-1}[M/N] = M.
    \]
    Furthermore, because $Q_{j-1} \subsetneq Q_j$ is a strict inclusion for each $1 \le j \le s$, their preimages must also satisfy strict inclusion:
    \[
        \pi^{-1}[Q_{j-1}] \subsetneq \pi^{-1}[Q_j].
    \]
    Concatenating the composition series of $N$ with the preimages of the composition series of $M/N$ gives the chain of submodules in $M$:
    \[
        0 = N_0 \subsetneq N_1 \subsetneq \cdots \subsetneq N_r = N = \pi^{-1}[Q_0] \subsetneq \pi^{-1}[Q_1] \subsetneq \cdots \subsetneq \pi^{-1}[Q_s] = M.
    \]
    To verify that this chain forms a composition series for $M$, we must check that every successive quotient is simple:
    \begin{enumerate}
        \item For $1 \le i \le r$, the quotient $N_i / N_{i-1}$ is simple by the assumption that $(N_i)_{i=0}^r$ is a composition series for $N$.
        \item For $1 \le j \le s$, define the map
        \[
            \tilde{\varphi}_j : \pi^{-1}[Q_j] \longrightarrow \frac{Q_j}{Q_{j-1}} \quad \text{by} \quad \tilde{\varphi}_j(m) = \pi(m) + Q_{j-1}.
        \]
        This is an $R$-module homomorphism, being the composition of the quotient map $\pi$ with the canonical projection $Q_j \to Q_j / Q_{j-1}$.
        
        \begin{itemize}
            \item \textbf{Surjectivity:} Take any element $x + Q_{j-1} \in Q_j / Q_{j-1}$ with $x \in Q_j$. Since $\pi$ is surjective, there exists $m \in M$ such that $\pi(m) = x$. By definition, $m \in \pi^{-1}[Q_j]$, and $\tilde{\varphi}_j(m) = x + Q_{j-1}$. Thus, $\tilde{\varphi}_j$ is surjective.
            \item \textbf{Kernel:} An element $m \in \pi^{-1}[Q_j]$ lies in $\ker \tilde{\varphi}_j$ if and only if $\pi(m) + Q_{j-1} = 0 \in Q_j / Q_{j-1}$. This is equivalent to $\pi(m) \in Q_{j-1}$, which means $m \in \pi^{-1}[Q_{j-1}]$. Thus, $\ker \tilde{\varphi}_j = \pi^{-1}[Q_{j-1}]$.
        \end{itemize}

        By the First Isomorphism Theorem, $\tilde{\varphi}_j$ induces an isomorphism of $R$-modules:
        \[
            \frac{\pi^{-1}[Q_j]}{\pi^{-1}[Q_{j-1}]} \cong \frac{Q_j}{Q_{j-1}}.
        \]
        Since $Q_j / Q_{j-1}$ is simple by the definition of the composition series for $M/N$, the factor module $\pi^{-1}[Q_j] / \pi^{-1}[Q_{j-1}]$ is also simple.
    \end{enumerate}

    Every factor in the spliced chain is simple, so the sequence of submodules
    \[
        0 = N_0 \subsetneq N_1 \subsetneq \cdots \subsetneq N_r = \pi^{-1}[Q_0] \subsetneq \pi^{-1}[Q_1] \subsetneq \cdots \subsetneq \pi^{-1}[Q_s] = M
    \]
    is a valid composition series for $M$.

    Finally, by the Jordan--Hölder Theorem, all composition series of $M$ have the same length $\ell(M)$, which is equal to the number of simple factors. Counting the factors in our constructed series gives:
    \[
        \ell(M) = r + s = \ell(N) + \ell(M/N).
    \]
\end{proof}

\begin{prob}
    Let $M = \mathbb Z/36\mathbb Z$ as a 
    $\mathbb Z$-module and let 
    $$
        N = \langle \overline 6 \rangle \leq M. 
    $$
    Determine $N$ then identify it up to isomorphism. 

    Find $\ell (N), \ \ell(M/N), \ \ell(M)$. 

    Finally check that 
    $$
        \ell(M) = \ell(N) + \ell(M/N).
    $$
\end{prob}
\begin{proof}
    Firstly we have 
    \[
        N = \langle \overline 6\rangle 
        = \{\overline 0, \overline 6, 
        \overline{12}, \ldots, \overline{30}\}, 
    \] 
    as this is an abelian group, we can use the 
    classification theorem for finitely generated abelian 
    groups to get 
    \[
        N \cong \mathbb Z_2 \times \mathbb Z_3 
        \cong \mathbb Z_6. 
    \]
    Now we want to find a maximal submodule of $M$. 
    Clearly 
    \[
        \langle \overline 2\rangle 
        = \{\overline 0, \overline 2, \ldots, 
        \overline{34}\}
    \]
    is such a submodule and it is isomorphic to 
    $\mathbb Z_{18}$ since it is a cyclic 
    abelian group of order $18$. 

    Lastly if we take $\langle \overline{18}\rangle$ we get 
    a composition chain 
    \[
        0 < \langle \overline{18} \rangle < 
        N < \langle \overline 2 \rangle < M. 
    \]
    So $\ell(M) = 4$. 
    Now 
    \[
        M/N = \frac{\mathbb Z/36\mathbb Z}{\mathbb Z_6}
        \cong \mathbb Z_6
    \]
    so $M/N \cong N$.  
    
    $\ell(M/N) = \ell(N) = 2$ since 
    \[
        0 < \langle \overline 3 \rangle < \langle \overline 
        2 \rangle < N = M/N
    \]
    is a valid chain. 

    Then $4 = \ell(M) = \ell(M/N) + \ell(N) = 2 + 2$ is 
    satisfied. 
\end{proof}

\begin{prob}
    Let $M$ be an $R$-module of finite length, and 
    let     
    \[
        f : M \to M 
    \]
    be an $R$-module homomorphism (endomorphism). 

    Prove that 
    $f$ is injective if and only if 
    $f$ is surjective. 

    You may use 
    \[
        M/\ker f \cong \im f 
    \]
    but no assumptions about exact sequences. 
\end{prob}
\begin{proof}
    $M$ has finite length so we can 
    find some composition series 
    \[
        0 = M_0 < M_1 < \cdots < M_n = M. 
    \]
    We are given that 
    \[
        M/\ker f \cong \im f. 
    \]
    
    Suppose $f$ is injective. Then $\ker f = \{0\}$. 
    Since 
    \[
        M/\ker f \cong \im f 
    \]
    we have 
    \[
        M = M/\ker f \cong \im f  
    \]
    and conversely $M = \im f$ would imply 
    \[
        M/\ker f \cong M, 
    \]
    but isomorphism doesn't necessarily imply equality.
    
    Suppose first $f$ is injective. Then 
    $\ker f = 0$ so 
    \[
        \ell(M) = \ell(0) + \ell(\im f) = \ell(\im f). 
    \]
    But $\im f \leq M.$ If $f$ were a proper submodule 
    it would have strictly less length. 
    Since their length are equal we have $\im f = M$. 
    So $f$ is surjective. 
    
    Conversely, suppose $f$ is surjective. Then 
    $\im f = M$ so we get 
    \[
        \ell(M) = \ell(\ker f) + \ell(M) 
    \]
    so $\ker f$ has length $0$, hence $\ker f = 0$. 
\end{proof}

\newpage 

\section{Exact sequences}

An exact sequence records, in one line, how 
kernels, images, submodules, quotients, injections, 
and surjections fit together. 

Gathmann introduces exact sequences as a 
unified graphical language for module constructions 
that would otherwise require a lot of repetitive 
prose. 

\begin{defn}
    Suppose we have $R$-modules and 
    $R$-module homomorphisms 
    \[
        M_1 \overset{\phi_1}{\to} M_2 
        \overset{\phi_2}{\to} M_3. 
    \]
    We say that the sequence is exact at $M_2$ 
    if 
    \[
        \im \phi_1 = \ker \phi_2. 
    \] 
    A longer sequence 
    \[
        M_1 \to M_2 \to \cdots \to M_n
    \]
    is exact if this equality holds at every 
    interior module. 
\end{defn}

So if 
\[
    A \overset{f}{\to} B \overset{g}{\to} C
\]
is exact at $B$, then 
\[
    \im f = \ker g. 
\]

In particular, 
\[
    g(f(a)) = 0
\]
for every $a \in A$, because 
\[
    \im f \subseteq \ker g. 
\]
But exactness says something stronger than 
merely $g \circ f = 0$. It says every element killed 
by $g$ must have come from $A$. 

\subsection{Exactness and the zero module}

The zero module $0$ is very useful in exact sequences 
as there is only one homomorphism to or from it. 

Consider 
\[
    0 \to M \overset{\varphi}{\to} N. 
\]
Exactness at $M$ says 
\[
    \im(0 \to M) = \ker \varphi. 
\]
But 
\[
    \im(0 \to M) = 0. 
\]
Therefore 
\[
    \ker \varphi = 0. 
\]
So $\varphi$ is injective. 

Thus 
\[
    0 \to M \overset{\varphi}{\to} N 
\]
is exact if and only if 
$\varphi$ is injective. 

Similarly, 
\[
    M \overset{\varphi}{\to} N \to 0
\]
is exact if and only if $\im \varphi = N$ 
which is that $\varphi$ is surjective. 

Thus exactness of  
\[
    0 \to M \overset{\varphi}{\to} N \to 0
\]
holds if and only if $\varphi$ is an isomorphism. 

\subsection{Short exact sequences}

The most important exact sequences are the 
short exact sequences 
\[
    0 \to M \overset{\varphi}{\to} N 
    \overset{\psi}{\to} P \to 0. 
\]
Exactness tells us three things: 
\begin{align*}
    \varphi \ \text{is injective}, \\ 
    \im \varphi = \ker \psi, 
\end{align*}
and 
\[
    \psi \ \text{is surjective}. 
\]

This has a rather nice interpretation. 

Because $\varphi$ is injective, we know 
\[
    M \subseteq N. 
\]
Because 
\[
    \ker \psi = \im \varphi \cong M, 
\]
the map $\psi$ kills exactly the copy of $M$ inside 
$N$. 

And because $\psi$ is surjective, 
\[
    P \cong N / \ker \psi. 
\]
Therefore 
\[
    P \cong N/M. 
\]
So conceptually, a short exact sequence 
\[
    0 \to M \to N \to P \to 0
\]
means $N$ contains a copy of $M$, and after 
quotienting $N$ by that copy, what remains is $P$. 

\subsection{The canonical short exact sequence}

Suppose $N \leq M$ is a submodule. 

Then there is always a short exact sequence 
\[
    0 \to N \overset{\iota}{\to} M 
    \overset{\pi}{\to} M/N \to 0. 
\]

Here $\iota(n) = n$ 
is the inclusin map and 
\[
    \pi(m) = m+N
\]
is the quotient map. 

The inclusion is injective, the quotient map 
is surjective and 
\[
    \ker \pi = N. 
\]

Since $\im(\iota) = N$ we have 
\[
    \im \iota = \ker \pi. 
\]

\subsection{A homomorphism produces 
exact sequences}

Take any module homomorphism $\varphi : M \to N$. 
We always have 
\[
    0 \to \ker \varphi \to M \overset{\varphi}{\to} 
    \im \varphi \to 0. 
\]
And also 
\[
    0 \to \im \varphi \to N \to N/\im \varphi \to 0. 
\]
Gluing them gives 
\[
    0 \to \ker \varphi \to M \overset{\varphi}{\to}
    N \to N/\im\varphi \to 0. 
\]

The last module 
\[
    N/\im \varphi
\]
is called the cokernel of $\varphi$: 
\[
    \operatorname{coker} \varphi := N/\im \varphi. 
\]

So another very useful way to package a map is 
\[
    0 \to \ker \varphi \to 
    M \overset{\varphi}{\to} N \to \operatorname{coker} 
    \to 0. 
\]
This tells us that the map has two kinds of failure. 

$\ker \varphi$ measures failure of injectivity, while 
$\operatorname{coker}\varphi$ measures failure of surjectivity.

\subsection{Splitting and gluing exact sequences}


Suppose 
\[
    M_1 \overset{\varphi_1}\to M_2 
    \overset{\varphi_2}\to M_3 
    \overset{\varphi_3}\to M_4
\]
is exact. 

Then at $M_3$, 
\[
    \im \varphi_2 = \ker \varphi_3.
\] 

Let $N := \im \varphi_2 = \ker \varphi_3$. 

Then we can split the sequence into 
\[
    M_1 \to M_2 \to N \to 0 
\]
and 
\[
    0 \to N \to M_3 \to M_4. 
\]

Conversely, those pieces can be glued back together. 

\subsection{Exactness and length}

If 
\[
    0 \to M \to N \to P \to 0 
\]
is exact and all modules have finite length, 
then 
\[
    \ell(N) = \ell(M) + \ell(P). 
\]

For a longer exact sequence 
\[
    0 \to M_1 \to M_2 \to \cdots \to M_n \to 0, 
\]
Gathmann derives the alternating-sum formula 
\[
    \sum_{i=1}^n (-1)^i\ell(M_i) = 0. 
\]

\subsection{Can the middle module be determined}

Suppose 
\[
    0 \to M \to N \to P \to 0. 
\]

A tempting guess is 
\[
    N \overset{\text{?}}{\cong} M \oplus P. 
\]
This is sadly not always true. 

There is ceirtanly always the obvious 
example 
\[
    0 \to M \to M \oplus P \to P \to 0, 
\]
where the first map is the inclusion 
$m \mapsto (m,0)$ and the second
is the projection $(m,p) \mapsto p$. 

Gathmann produces a counter-example to 
this idea. 

\subsection{Split exact sequences}

A short exact sequence 
\[
    0 \to M \overset{\varphi}\to N \overset{\psi}\to P 
    \to 0 
\]
is called split exact when 
\[
    N \cong M \oplus P
\]
in a way compatible witht the maps. 

Gathmann's Splitting Lemma says the 
following conditions 
are equivalent: 
\[
    N \cong M \oplus P,
\]
$\varphi$ has a left inverse 
\[
    \alpha : N \to M, \ \alpha \circ \varphi = \operatorname{id}_M,
\]
or $\psi$ has a right inverse 
\[
    \beta: P \to N, \ \psi \circ \beta = \operatorname{id}_P. 
\]
 
The map $\beta : P \to N$ is often called a section. 

It means that we can choose, in a homomorphic way, 
one representative in $N$ for every element of $P$. 

Then $N$ decomposes into 
the kernel part of $M$ plus a copy of $P$.  

For vector spaces, an exact sequence 
always splits. 

\subsection{Commutative diagrams}

\begin{defn}
    A commutative diagram is a diagram of modules 
    and maps in which every path from 
    one objects to another gives the same map. 
\end{defn}

For example, 
\[
\begin{tikzcd}
A \arrow[r, "f"] \arrow[d, "g"'] & B \arrow[d, "h"] \\
C \arrow[r, "k"']                 & D
\end{tikzcd}
\]
is commutative when 
\[
    h \circ f = k \circ g. 
\]

Suppose we have two exact rows and vertical 
homomorphisms 

\[
\begin{tikzcd}
M \arrow[r, ""] \arrow[d, "\alpha"] & N \arrow[d, "\beta"] 
\arrow[r, ""] & P \arrow[d, "\gamma"] \arrow[r,""] & 0 \\
0 \arrow[r, ""] & M' \arrow[r, ""] & N' \arrow[r, ""] 
& P' 
\end{tikzcd}
\]
with the squares commutative. 

The Snake Lemma manufactures a new exact sequence 
involving the kernels and cokernels of the 
vertical maps: 
\[
    \ker \alpha \to \ker \beta 
    \to \ker \gamma \overset{\delta}\to M'/\im\alpha 
    \to N'/\im\beta \to P'\im\gamma. 
\]

In the notation introduced earlier we write 
\[
    \ker \alpha \to \ker \beta \to \ker \gamma 
    \overset{\delta}\to \coker \alpha 
    \to \coker \beta \to \coker \gamma. 
\]
The map $\delta : \ker \gamma \to \coker \alpha$ is 
called the connecting homomorphism. 

Conceptually, if you have 
two exact rows linked by vertical maps, 
the Snake Lemma relates the failures of those 
vertical maps to be injective and surjective. 

Kernels measure failure of injectivity, cokernels 
measure the failure of surjectivity. The Snake Lemma 
ties them nicely together. 

Gathmann constructs $\gamma$ by doing a diagram
chase: life an element backwards through one map, 
move vertically and life backwards again. 

\subsection{Some problems to see if I understand what the hell is going on}

\begin{prob}
    Let 
    \[
        A \overset{f}\to B \overset{g}\to C 
    \]
    be a sequence of $R$-modules. 
    Suppose 
    \[
        \im f = \ker g. 
    \]
    Answer the following: 
    \begin{enumerate}
        \item At which module is the sequence exact? 
        \item Must $f$ be injective? 
        \item Must $g$ be surjective? 
        \item Prove that 
            \[
                g \circ f = 0. 
            \]
    \end{enumerate}
\end{prob}
\begin{proof}[Solution and proof sketch]
    The sequence is exact at 
    $B$. 

    $f$ needs not be injective and $g$ 
    needs not be surjective. 

    We have $\im f = \ker g$. 
    Then $g(x) = 0$ for all $x = f(a) \in \im f$. 
    Then $g \circ f = 0$. 
\end{proof}

\begin{prob}
    Let 
    \[
        0 \to M \overset{\varphi}\to N 
    \]
    be exact. 

    Prove directly from the definition that 
    $\varphi$ is injective. 

    Then let 
    \[
        M \overset{\psi}\to N \to 0 
    \]
    be exact and prove that $\psi$ is surjective.
    
    Finally explain why 
    \[
        0 \to M \overset{\alpha}\to N \to 0 
    \]
    is exact if and only if $\alpha$ is an isomorphism. 
\end{prob}
\begin{proof}[Proof sketch]
    There exists precisely 
    one homomorphism from $0 \to M$, that 
    being the $0$-homomorphism. 
    Thus $\im(0 \to M) = \{0\}$. 
    Then, if the sequence is exact, then 
    $\{0\} = \im(0 \to M) = \ker \varphi$ 
    so $\varphi$ is injective. 

    Likewise, there exists one and only one 
    homomorphism from $N$ to $0$, that being 
    the $0$-homomorphism. 
    Thus, if $\im \psi = \ker (N \to 0) = N$ so 
    $\psi$ is surjective. 
    
    Summing up the previous two it is easy to see 
    that we require 
    $\ker \alpha = \{0\}$ and $\im \alpha = N$ 
    so it is bijective and thus an isomorphism. 

    Suppose now that $\alpha$ is an isomorphism. 
    Then $\ker \alpha = \{0\}$ and 
    $\im \alpha = N$ which, together with the fact 
    that only the $0$-homomorphism is valid from 
    $0 \to M$ and $N \to 0$ gives the 
    equivalence we want. 
\end{proof}

\begin{prob}
    Let $N \leq M$ be a submodule. 
    Consider 
    \[
        0 \to N \overset{\iota}\to M 
        \overset\pi\to M/N \to 0, 
    \]
    where $\iota(n) = n$ and $\pi(m) = m+N$. 
    Prove this sequence is exact. 
\end{prob}
\begin{proof}
    Trivially, 
    $\im (0 \to N) = \{0\}$ and 
    $\ker (M/N \to 0) = M/N$. 

    If $\ker \iota \supsetneq \{0\}$ then 
    there would be nonzero $n \in M$ 
    such that $n = 0$ which would imply $M = N = 0$, 
    which in turn would imply $\ker \iota = \{0\}$. 
    (This was a really dumb argument, but it is obvious 
    that $\ker \iota = 0$). 

    Next, $\pi$ is surjective ($\im \pi = M/N$) 
    as we have seen before.   

    It remains to show that 
    $\im \iota = \ker \pi$. 

    Suppose $\pi m = 0$. Then $\overline m + N 
    = 0 + N$ so 
    $m \in N$ so $\ker \iota = N$. 
    
    $\im \iota = N$ is trivial to see. 
\end{proof}

\begin{prob}
    Consider 
    \[
        0 \to \mathbb Z/2\mathbb Z 
        \overset\varphi\to \mathbb Z/4\mathbb Z 
        \overset\psi\to \mathbb Z/2\mathbb Z \to 0,
    \]
    where $\varphi(\overline 1) = \overline 2$ and 
    $\psi$ is reduction modulo $2$. 
    Prove it is exact. 

    Decide whether it is split exact. 
\end{prob}
\begin{proof}
    We have 
    $\ker \varphi = \overline 0 = \im(0 \to \mathbb Z/
    2\mathbb Z)$.  Next, 
    $\im \varphi = \{\overline 0, \overline 2\} = \ker \psi$ 
    since $\psi$ is reduction modulo $2$. 
    Lastly $\im \psi = \mathbb Z/2\mathbb Z = 
    \ker(\mathbb Z/2\mathbb Z \to 0)$. 
    
    This sequence is not split 
    exact because $\mathbb Z/2\mathbb Z 
    \oplus \mathbb Z/2\mathbb Z 
    \not\cong \mathbb Z/4\mathbb Z$. (The 
    latter has an element of order $4$, the former 
    does not). 
\end{proof}

\begin{prob}
    Let $I,J \triangleleft R$ be ideals. 
    Prove there is a short exact sequence 
    \[
        0 \to I \cap J 
        \to I \oplus J \to I+J \to 0. 
    \]
\end{prob}
\begin{proof}
    Consider $\varphi : I \cap J \to I \oplus J$
    and $\psi : I \oplus J \to I + J$ by 
    \[
        \varphi(x) = (x, -x), 
    \] 
    and 
    \[
        \psi(x,y) = x+y. 
    \]

    $\im(0 \to I \cap J) = \{0\}$ and 
    $\ker \varphi = \{0\}$. 
    Also $\im \varphi = \{(n,n) \bigm| n \in I \ \text{and} \ 
    n \in J\}$. Then it is easy to see that 
    $\ker \psi = \im \varphi$. The final 
    correspondence is similarly easy. 
\end{proof}

\begin{prob}
    Let $R = \mathbb Z$ with ideal 
    $I = (5)$. Let $M = \mathbb Z_{12}$. 
    Show that $M = IM$. Can you find an 
    element $a \in I$ such that 
    $am = m$ for every $m \in M$. 
\end{prob}
\begin{proof}[Solution and proof]
    Recall that $IM$ is the 
    set of all finite sums $\sum a_im_i$ where 
    $a_i \in I, \ m_i \in M$. 

    First we find an 
    $a \in I$ so that $am = m$. 
    Notice that $a:= 25 \in I$ is 
    $1$ modulo $12$ so 
    \[
        am \equiv 1 \cdot m = m. 
    \]

    This means that for every 
    $m \in M$ we have 
    \[
        m = 25m \in IM  
    \]
    so $M \subseteq IM$. 
    The reverse inclusion is automatic from 
    the definition of an $R$-module. 
\end{proof}

\section{Tensor Products}
We have covered machinery for modules and linear 
maps, but many natural maps are not linear in 
one variable. Rather, they are bilinear. 

For instance if $M,N,P$ are $R$-modules, a map 
\[
    \alpha : M \times N \to P
\]
is $R$-bilinear if, for fixed $n$, the map 
\[
    m \mapsto \alpha(m,n)
\]
is linear, and for fixed $m$, the 
map 
\[
    n \mapsto \alpha(m,n)
\]
is linear.  

We want to build a new module 
$M \otimes_R N$ such that studying 
bilinear maps out of $M \times N$ is 
equivalent to studying ordinary linear maps 
out of $M \otimes_R N$. 

More precisely, we want 
\[
    \{\text{bilinear maps} \ M \times N \to P\} 
    \leftrightarrow 
    \{\text{linear maps} \ M \otimes_R N \to P\}.
\]

\subsection{Motivation from free modules}
Suppose first that $M = R^m$ and $N = R^n$ 
with bases 
\[
    b_1,\ldots,b_m
\]
and 
\[
    c_1,\ldots,c_n.
\]

A bilinear map 
\[
    \alpha : M \times N \to P
\]
is determined by the $mn$ values 
\[
    \alpha(b_i,c_j). 
\]
Indeed 
\[
    \alpha\left(\sum_i \lambda_i b_i, \sum_j \mu_j c_j\right) 
    = \sum_{i,j}\lambda_i \mu_j \alpha (b_i,c_j).
\]
We imagine constructing a free module with basis 
symbols $b_i \otimes c_j$.
Call it $F$. Then a linear map 
\[
    F \to P
\]
is also determined by choosing arbitrary values for those 
$mn$ basis elements. 
Thus 
\[
    \alpha(b_i,c_j) \leftrightarrow 
    \varphi(b_i \otimes c_j). 
\]
So for free modules it should be natural that 
\[
    R^m \otimes_R R^n \cong R^{mn}. 
\]

\subsection{The universal property}
\begin{defn}[Tensor product]
    A \textbf{tensor product} of $M$ and $N$ 
    over $R$ consists of an $R$-module 
    \[
        M \otimes_R N
    \]
    together with a bilinear map 
    \[
        \tau : M \times N \to M \otimes_R N, 
        \ (m,n) \mapsto m \otimes N, 
    \]
    such that every bilinear map 
    \[
        \alpha : M \times N \to P
    \]
    factors uniquely through a linear map 
    \[
        \varphi : M \otimes_R N \to P. 
    \]
    Diagramatically, 
    \[
    \begin{tikzcd}
    M\times N \arrow[r, "\alpha"] 
    \arrow[dr, "\tau"] & P \\
        & M \otimes_R N \arrow[u, "\varphi"]
    \end{tikzcd}
    \]
    with 
    \[
        \alpha = \varphi \circ \tau. 
    \]
    This is called the universal property. 
\end{defn}
Once we know what $\alpha$ should do, the universal 
product gives us a unique linear map from the 
tensor product. 

This definition tells us what it does rather 
than giving a construction. 
This is sufficient because the definition uniquely 
determines the tensor product up to isomorphism. 
So if two modules $T_1, \ T_2$ perform this job then there 
is a unique isomorphism $T_1 \cong T_2$ comptaible 
with their tensor maps. 

The symbol $m \otimes n$ is not an ordered pair.  
It denotes an element of the new module 
$M \otimes_R N$. 
The defining bilinearity gives the relations 
\[
    (m_1 + m_2) \otimes n = m_1 \otimes n + 
    m_2 \otimes n  
\]
and 
\[
    m \otimes (n_1 + n_2) 
    = m \otimes n_1 + m \otimes n_2,  
\]
and 
\[
    (am) \otimes n = a(m\otimes n) = m\otimes (an).  
\]

\subsection{Pure tensors vs. global tensors} 
An element specifically of the form $m \otimes n$ 
is called a pure tensor. 
But not every element of $M \otimes_R N$ is of this form. 
A general tensor is a finite sum 
\[
    \sum_{i=1}^k a_i(m_i \otimes n_i). 
\]
Pure tensors generate the tensor product, but 
don't necessarily form a basis, and representations 
are genrally not unique. 

\subsection{The construction}
Start with a large free $R$-module $F$ whose basis 
is indexed by all pairs 
\[
    (m,n) \in M \times N. 
\]
At this stage,
\[
    (m_1 + m_2, n)
\]
and 
\[
    (m_1, n) + (m_2,n)
\]
are totally unrelated formal expressions.
We want them to become equal, because tensoring 
should be bilinear. 
Take a submodule $G$ generated by all expressions forcing 
bilinearity 
\[
    (m_1+m_2, n) - (m_1,n) - (m_2,n), 
\]
\[
    (am,n) - a(m,n), 
\]
\[
    (m,n_1+n_2) - (m,n_1) - (m,n_2), 
\]
\[
    (m,an) - a(m,n). 
\]
Then define 
\[
    M \otimes_R N := F/G. 
\]

\subsection{Three isomorphisms}
First, 
\[
    M \otimes_R N \cong N \otimes_R M. 
\]
The map is $m \otimes n \mapsto n \otimes m$. 

Second, 
\[
    M \otimes_R R \cong M. 
\]
The map is $m \otimes a \mapsto am$. 
The inverse map is $m \mapsto m \otimes 1$. 

Third, tensor products distribute over direct sums 
\[
    (M \oplus N) \otimes_R P 
    \cong (M \otimes_R P) \oplus (N \otimes_R P). 
\]

Also, tensor products are associative up to natural 
isomorphism 
\[
    (M \otimes N) \otimes P = 
    M \otimes (N \otimes P)
\]
so we usually just write 
\[
    M \otimes N \otimes P. 
\]

\subsection{Free modules}
If $M \cong R^m$ and $N \cong R^n$ 
then 
\[
    M \otimes N \cong R^{mn}.
\]
If $b_1,\ldots,b_m$, $c_1,\ldots,c_n$ are bases, then 
\[
    b_i \otimes c_j
\]
forms a basis of the tensor product. 
So over a field, 
\[
    \dim(V \otimes W) = \dim V \cdot \dim W. 
\]

\subsection{Vanishing}
Suppose $I,J \triangleleft R$ are coprime ideals, meaning 
\[
    I + J = R. 
\]
Then 
\[
    R/I \otimes_R R/J = 0. 
\]
Choose $a \in I, b \in J$ with $a+b = 1$. 
Then, for any pure tensor 
\[
    \overline{r}\otimes \overline s 
    = (a+b)(\overline r \otimes \overline s) 
    = a \overline r \otimes \overline s + 
    \overline r \otimes b \overline s. 
\]
But $a \overline r = 0$ in $R/I$ and 
$b\overline s = 0$ in $R/J$.
Therefore $\overline r \otimes \overline s = 0$. 

So for distinct primes $p,q$ we have 
\[
    \mathbb Z/p\mathbb Z \otimes \mathbb Z/q\mathbb Z =0. 
\]

Sidenote: 

In general
\[
    M' \subseteq M \not\Rightarrow 
    M' \otimes N \subseteq M \otimes N. 
\]

\subsection{Tensor product of homomorphisms}
Suppose $\varphi: M \to N$ 
and $\varphi' : M' \to N'$ are module homomorphisms. 
Then there is a natural homomorphism 
\[
    \varphi \otimes \varphi' : M \otimes M' \to 
    N \otimes N' 
\]
defined on pure tensors by $(\varphi \otimes \varphi')(m \otimes m') 
= \varphi(m) \otimes \varphi(m')$. 

This is well-defined because the map 
\[
    (m,m') \mapsto \varphi(m) \otimes \varphi(m')
\]
is bilinear, so the universal property gives the 
corresponding linear map. 

\subsection{Extension of scalars}
Suppose $R'$ is an $R$-algebra and $M$ is an 
$R$-module. 

We can turn $M$ into an $R'$-module by definining 
\[
    M_{R'} := M \otimes_R R'. 
\]
This is called an extension of scalars. 
The new scalar multiplication is 
\[
    a \cdot (m \otimes s) = m \otimes (as), \ 
    a,s \in R'. 
\]
As a motivating example, consider a real vector space $V$. 
To allow complex coefficients, define 
\[
    V_{\mathbb C} = V \otimes_{\mathbb R} \mathbb C. 
\]
This is the complexification of $V$. 
If $\dim_{\mathbb R}V = n$ then 
\[
    V_{\mathbb C} \cong \mathbb C^n. 
\]

\subsection{Tensor products of algebras}
If $R_1,R_2$ are not merely $R$-modules, but $R$-algebras, 
then 
\[
    R_1 \otimes R_2
\]
itself is an $R$-algebra. 

Define multiplication by pure tensors as 
\[
    (s \otimes t)(s'\otimes t') = ss' \otimes tt'. 
\]
Bilinearity extends to arbitrary tensors. 
A nice example is 
\[
    R[x,y] \cong R[x] \otimes_R R[y]. 
\]
The map is 
\[
    f(x)\otimes g(x) = f(x)g(x). 
\]
In particular, 
\[
    x^i \otimes y^j = x^iy^j. 
\]

\subsection{The Hom-tensor relation}
There is a natural isomorphism 
\[
    \operatorname{Hom}_R(M,\operatorname{Hom}_R(N,P)) 
    \cong \operatorname{Hom}_R(M \otimes_R N, P). 
\]
Suppose $\alpha : M \to \operatorname{Hom}(N,P)$. 
For each $m, \ \alpha(m)$ is itself a map $N \to P$, so we 
can define 
\[
    \beta(m\otimes n) = \alpha(m)(n). 
\]
Conversely, given 
\[
    \beta : M \otimes N \to P,
\]
define 
\[
    \alpha(m)(n) = \beta(m\otimes n). 
\]

\subsection{Tensor products are right exact}
Suppose 
\[
    M_1 \overset{\varphi_1}\to M_2 
    \overset{\varphi_2}\to M_3 \to 0
\]
is exact. 
Tensor everything with $N$, 
\[
    M_1 \otimes N 
    \overset{\varphi_1 \otimes \operatorname{id}} M_2 \otimes N 
    \overset{\varphi_2 \otimes \operatorname{id}} M_3 \otimes N \to 0. 
\]
Then this sequence is exact. 

Interestingly, tensor products are not left exact. 


\end{document}
